\phantomsection\label{conj:master-infinite-generator}%
\phantomsection\label{constr:platonic-package}%
\phantomsection\label{conv:bar-coalgebra-identity}%
\phantomsection\label{conv:hms-levels}%
\phantomsection\label{conv:regime-tags}%
\phantomsection\label{cor:bar-is-dgcoalg}%
\phantomsection\label{cor:shadow-extraction}%
\phantomsection\label{def:chiral-ass-operad}%
\phantomsection\label{def:cyclically-admissible}%
\phantomsection\label{def:shadow-depth-classification}%
\phantomsection\label{def:shadow-postnikov-tower}%
\phantomsection\label{prop:independent-sum-factorization}%
\phantomsection\label{prop:pbw-universality}%
\phantomsection\label{rem:five-pieces}%
\phantomsection\label{rem:yangian-three-theorems}%
\phantomsection\label{sec:bar-nilpotency-nine-terms-complete}%
\phantomsection\label{sec:heisenberg-genus-expansion-eisenstein}%
\phantomsection\label{sec:k3-sigma-model}%
\phantomsection\label{thm:arnold-relations}%
\phantomsection\label{thm:bar-cobar-adjunction}%
\phantomsection\label{thm:bar-cobar-verdier}%
\phantomsection\label{thm:bar-nilpotency-complete}%
\phantomsection\label{thm:e1-chiral-koszul-duality}%
\phantomsection\label{thm:general-hs-sewing}%
\phantomsection\label{thm:genus-universality}%
\phantomsection\label{thm:heisenberg-genus-one-complete}%
\phantomsection\label{thm:heisenberg-genus2-obstruction}%
\phantomsection\label{thm:km-quantum-groups}%
\phantomsection\label{thm:lattice-voa-bar}%
\phantomsection\label{thm:lattice:curvature-braiding-orthogonal}%
\phantomsection\label{thm:modular-characteristic}%
\phantomsection\label{thm:planted-forest-tropicalization}%
\phantomsection\label{thm:riccati-algebraicity}%
\phantomsection\label{thm:shadow-connection}%
\phantomsection\label{thm:yangian-bar-rtt}%
\phantomsection\label{thm:yangian-koszul-dual}%
\phantomsection\label{subsec:DDCA-ainfty}%
\phantomsection\label{thm:twisted-M2-primitive-package}%
\phantomsection\label{eq:DDCA-central-term}%
\phantomsection\label{eq:M2-r-matrix}%
\phantomsection\label{thm:DDCA-m2}%
\phantomsection\label{eq:DDCA-bracket}%
\phantomsection\label{subsec:M2-holography}%

\chapter{Toroidal and elliptic algebras}
\label{chap:toroidal-elliptic}

The chiral algebras of Vols~I--II each depend on a single deformation
parameter: the level~$k$, the central charge~$c$, or the spectral
parameter~$u$. This chapter enters the regime where one parameter is
not enough. The toroidal algebra
$U_{q,t}(\hat{\hat{\fg}})$ carries two independent parameters $(q,t)$,
the elliptic quantum group $E_{q,p}(\fg)$ replaces the rational
$R$-matrix by a doubly-periodic one, and the double affine Hecke
algebra (DAHA) intertwines spectral and modular directions. In these algebras the modular parameter~$\tau$ and the spectral
parameter~$u$ are coupled through the Fay trisecant identity and the
quasi-periodicity of theta functions.

These algebras are shadows of the canonical stage~$1$ factorisation algebra $\PhiFA_d(\cC) \in \EdHolFA(X)$ of Theorem~\ref{thm:phi-two-stage-factorisation}, read through the specialisation $\SpCh_{\Sigma_{d-1}, C}$ at the elliptic chart of its CY$_d$ base. For a framed CY$_3$ such as $K3 \times E$, the stage-$1$ holomorphic factorisation algebra on the threefold specialises along $(\Sigma_2,C)$ to an $\Eone$-chiral algebra on the curve $C = E_\tau$; the $\Etwo$-braiding attached to its $R$-matrix belongs to the centre side $\cZ(\Rep^{\Eone})$. On a torus-fibred chart the same specialisation mechanism supplies the expected toroidal current algebra. The canonical lane carries the higher-dimensional factorisation structure on the surface or threefold; the shadow lane lives on the elliptic reference curve.

The bar-cobar framework meets this two-parameter world in two ways.
Direction~(a) keeps the base a single curve, now of genus~$1$: the
propagator becomes elliptic, the Arnold relation fails by a
quasi-modular defect (Theorem~\ref{thm:elliptic-arnold-fails}) whose
repair is Felder's dynamical parameter, and the bar differential
acquires Eisenstein-series corrections indexed by~$\tau$. Direction~(b)
seeks genuine surface-factorization objects on $\bC^* \times \bC^*$,
requiring a two-dimensional factorization framework beyond the
curve-chiral hierarchy of Vols~I--II. Only direction~(a) is developed
here.

At genus~$1$ the three-term calculus of the Arnold relation survives
only in quartic form: the Fay trisecant identity
(Theorem~\ref{thm:fay}) relates products of four theta functions, and
its kernel form carries the dynamical variable inside the identity
itself. The bilinear three-term relation holds only after the
two-step degeneration $E_\tau \to \bC^*$ (as $\tau \to i\infty$,
the elliptic curve degenerates to a cylinder) followed by
$\bC^* \to \bC$ (as $q \to 1$, the cylinder opens to the affine line).

\begin{remark}[Curve and surface factorisation]
\label{rem:toroidal-two-tracks}
The curve construction keeps the base an algebraic curve and constructs
an elliptic $\Eone$-chiral realisation on $E_\tau$. The surface
construction asks for a genuine factorisation object on
$\bC^*\times\bC^*$. The results below concern the curve construction;
the surface construction remains a separate existence problem.
\end{remark}

\begin{remark}[The generalization principle: Arnold to Fay]
\label{rem:arnold-fay-generalization}
\index{Arnold relation!Fay generalization}
The Arnold relation (Vol~I, Theorem~\ref{thm:arnold-relations}) governs the genus-$0$ bar differential on rational curves. For curve-based $\Eone$-chiral realizations it has two generalisation axes. The \emph{modular axis}: via clutching morphisms and the shadow obstruction tower $\Theta_A^{\leq r}$ (Vol~I, Definition~\ref{def:shadow-postnikov-tower}), one passes to higher-genus stable curves, recovering the higher-genus bar complex of Vol~I (Vol~I, Remark~\ref{rem:five-pieces}). The \emph{elliptic-propagator axis}: replacing the additive logarithmic propagator $\omega = dz/(z-w)$ by the doubly-periodic propagator built from $\theta_1(z|\tau)$ makes the Arnold relation fail by the quasi-modular defect of Theorem~\ref{thm:elliptic-arnold-fails}; the failure is absorbed either as curvature (abelian targets) or by Felder's dynamical parameter (Corollary~\ref{cor:elliptic-bar-dynamical}), whose coherence rests on the quartic Fay trisecant identity (Theorem~\ref{thm:fay}). This enters the toroidal regime at fixed genus~$1$. The two axes meet in the higher-genus elliptic bar complex of Remark~\ref{rem:toroidal-two-tracks}.
\end{remark}

\begin{remark}[Rational, trigonometric, elliptic: curve geometry]
\label{rem:rational-trig-elliptic}
\index{quantum group hierarchy!curve geometry}
\index{R-matrix!rational, trigonometric, elliptic}
The three quantum group families arise from the geometry of the underlying curve:
\begin{center}
\small
\begin{tabular}{llll}
\textbf{$R$-matrix type} & \textbf{Quantum group} &
 \textbf{Spectral parameter} & \textbf{Curve} \\ \hline
Rational
 & Yangian $Y(\fg)$
 & $u = z_1 - z_2$
 & $\mathbb{P}^1$ \\[2pt]
Trigonometric
 & Quantum loop $U_q(L\fg)$
 & $q = e^{2\pi i(z_1 - z_2)/\beta}$
 & $\mathbb{C}^*$ \\[2pt]
Elliptic
 & Elliptic QG $E_{q,p}(\fg)$
 & $u \in E_\tau$
 & $E_\tau$
\end{tabular}
\end{center}
The passage rational $\to$ trigonometric $\to$ elliptic parallels additive $\to$ multiplicative $\to$ modular spectral parameter.
\end{remark}

\begin{remark}[Elliptic Hall algebra, spherical DAHA at genus~\texorpdfstring{$1$}{1}, quantum-toroidal \texorpdfstring{$(q,t)$}{(q,t)}]
\label{rem:elliptic-hall-spherical-daha-genus-1}
\index{elliptic Hall algebra}
\index{spherical DAHA!genus~$1$}
\index{quantum toroidal algebra!$(q,t)$ extension}
The $\mathfrak{gl}_1$ elliptic regime admits three presentations after the standard parameter identifications and completions. The \emph{elliptic Hall algebra} $E_{q,t}$ of Burban--Schiffmann (Duke Math.~J.~161 (2012), 1171--1231) is the shuffle/Hall positive half with structure function $\xi(u) = \theta(u+\gamma)/\theta(u)$; the \emph{spherical double affine Hecke algebra} $\mathcal{SH}(E_\tau)$ of Cherednik is its spherical polynomial-action presentation; the \emph{quantum toroidal algebra} $U_{q,t}(\hat{\hat{\mathfrak{gl}}}_1)$ of Ding--Iohara--Miki~\cite{DI97,Miki07} (Definition~\ref{def:toroidal}) is the Drinfeld-current presentation after adjoining the Cartan currents and completing. The identification
\[
 E_{q,t} \;=\; \mathcal{SH}(E_\tau) \;=\; U_{q,t}(\hat{\hat{\mathfrak{gl}}}_1)
\]
(Schiffmann--Vasserot~\cite{SV13}; Feigin--Tsymbaliuk~2011) is read in this completed $\mathfrak{gl}_1$ sense. For simply-laced $\mathfrak{g}$ beyond $\mathfrak{gl}_1$, Feigin--Jimbo--Miwa--Mukhin construct quantum-toroidal shuffle realizations of the positive half with a $\mathfrak{g}$-dependent kernel. A full compact Hall or PBW statement requires the negative half, Cartan part, Hopf pairing, topological completion, and compatibility of the induced braiding with the centre side $\cZ(\Rep^{\Eone})$. In every completed presentation the two parameters $(q,t)$ couple through the $\SL_2(\Z)$-Miki action of Remark~\ref{rem:miki-automorphism}; the Koszul involution $(q,t) \mapsto (q^{-1}, t^{-1})$ of Conjecture~\ref{conj:toroidal-koszul-dual} commutes with this action.
\end{remark}

\begin{construction}[Elliptic shadow data]
\label{constr:elliptic-shadow-data}
\index{elliptic quantum group!shadow data}
Toroidal and elliptic quantum group algebras carry shadow data in the elliptic regime:
\begin{enumerate}[label=\textup{(\roman*)}]
\item The prime form $E(z,w;\tau) = \theta_1(z-w;\tau)/\theta_1'(0;\tau)$ replaces the rational difference $z-w$; the propagator is its logarithmic differential $\eta^{\mathrm{ell}} = d\log\theta_1(z-w|\tau)$, replacing $dz/(z-w)$.
\item The genus-$1$ MC equation $D^{(1)}\Theta + \tfrac{1}{2}[\Theta^{(0)}, \Theta^{(0)}] + \kappa_{\mathrm{ch}}\,\omega_1 = 0$ carries a curvature term sourced by the quasi-modular defect $\pi^2 E_2(\tau)$ of the naive Arnold relation (Theorem~\ref{thm:elliptic-arnold-fails}, Computation~\ref{comp:ell-curvature}).
\item The naive elliptic Arnold sum $\eta_{12}\wedge\eta_{23} + \eta_{23}\wedge\eta_{31} + \eta_{31}\wedge\eta_{12}$ does \emph{not} vanish (Theorem~\ref{thm:elliptic-arnold-fails}); the CYBE is replaced by Felder's dynamical Yang--Baxter equation (Definition~\ref{def:dybe}), whose scalar entries reduce to the quartic Fay identity (Theorem~\ref{thm:fay}).
\end{enumerate}
\end{construction}

\section{Double affine setup}
\label{sec:double-affine}

Affinizing $\fg$ once produces $\widehat{\fg}$ with one spectral parameter $z$
and one level $k$. A Calabi--Yau threefold resolved by a single curve carries
only this one loop direction; a threefold resolved by a \emph{surface}
$\mathbb{C}^* \times \mathbb{C}^*$ carries two, and the algebra attached to
it must be affinized twice. The second loop obstructs naive iteration: the
horizontal and vertical currents cannot both be free because the
Calabi--Yau constraint $q_1 q_2 q_3 = 1$ forces the product of the three
directional weights to be trivial. Two independent parameters $(q,t)$
survive, with $t = q^k$ encoding the second level. The toroidal algebra
$U_{q,t}(\hat{\hat{\fg}})$ is the resulting double-loop structure.

\subsection{Definition}

\begin{definition}[Toroidal algebra]\label{def:toroidal}
\index{toroidal algebra|textbf}
\index{double affine algebra|textbf}
The \emph{toroidal algebra} (or double affine algebra) $U_{q,t}(\hat{\hat{\mathfrak{g}}})$
is generated by:
\begin{itemize}
\item Horizontal currents $x^\pm_{i,n}$, $\psi^\pm_{i,n}$ (affine in one direction)
\item Vertical currents $e_i(z)$, $f_i(z)$, $\psi_i(z)$ (affine in the other direction)
\end{itemize}
subject to:
\begin{enumerate}
\item \emph{Drinfeld--type relations} in each direction:
$[h_{i,m}, x^\pm_{j,n}] = \pm a_{ij} x^\pm_{j,m+n}$, \quad
$[x^+_{i,m}, x^-_{j,n}] = \delta_{ij} (\psi^+_{i,m+n} - \psi^-_{i,m+n})/(q_i - q_i^{-1})$
\item \emph{Cross relations} coupling horizontal and vertical:
$\psi_i(z) e_j(w) = g_{ij}(z/w) \, e_j(w) \psi_i(z)$
where $g_{ij}(u) = (u - q^{a_{ij}})/(u - q^{-a_{ij}})$
\item \emph{Serre relations} in both directions
\end{enumerate}
with two independent parameters $q$ and $t = q^k$. This is the finite-current
part of the presentation. The complete topological algebra also contains
the central and Cartan data, the Serre completions, and the Hopf structure
used to form doubles; see Ginzburg--Kapranov--Vasserot for the full
presentation.
\end{definition}

\begin{proposition}[Toroidal OPE]\label{prop:toroidal-ope}
In the standard current completion and Ding--Iohara normalisation, the
singular part of the OPE for toroidal currents involves both parameters:
\[
e_i(z) f_j(w) = \frac{\delta_{ij}}{(1 - q)(1 - t^{-1})} \cdot
\frac{\psi_i^+(w) - \psi_i^-(w)}{z - w} + \text{rational terms in } q, t
\]
\end{proposition}

\begin{proof}
We derive the OPE from the defining relations of Definition~\ref{def:toroidal}.

\emph{Arrow~1: Mode commutator.}
The Drinfeld-type relation (1) gives:
\[
[e_{i,m}, f_{j,n}] = \delta_{ij} \frac{\psi^+_{i,m+n} - \psi^-_{i,m+n}}{q_i - q_i^{-1}}.
\]
Form the generating functions $e_i(z) = \sum_m e_{i,m} z^{-m-1}$,
$f_j(w) = \sum_n f_{j,n} w^{-n-1}$, and
$\psi_i^\pm(w) = \sum_n \psi^\pm_{i,n} w^{-n}$.

\emph{Arrow~2: Extracting the OPE.}
Contracting the commutator against $z^{-m-1} w^{-n-1}$ and summing:
\[
e_i(z) f_j(w) - {:}e_i(z) f_j(w){:}
= \delta_{ij} \sum_{m} \frac{\psi^+_{i,m}(w) - \psi^-_{i,m}(w)}{q_i - q_i^{-1}} \cdot \frac{w^{-m}}{z - w}.
\]
For the toroidal algebra with parameters $q$ and $t = q^k$,
we have $q_i = q^{d_i}$ where $d_i = (\alpha_i, \alpha_i)/2$.
The singular part of the OPE is the residue at $z = w$:
\[
e_i(z) f_j(w) \sim \frac{\delta_{ij}}{q_i - q_i^{-1}} \cdot
\frac{\psi_i^+(w) - \psi_i^-(w)}{z - w}.
\]

\emph{Arrow~3: Two-parameter structure.}
The cross relation (2) introduces the second parameter.
From $\psi_i(z) e_j(w) = g_{ij}(z/w) \, e_j(w) \psi_i(z)$
with $g_{ij}(u) = (u - q^{a_{ij}})/(u - q^{-a_{ij}})$,
the normal ordering of $\psi$ against $e, f$ currents produces
corrections rational in both $q$ and $t$. Specifically, expanding
$1/(q_i - q_i^{-1}) = 1/((q^{d_i} - q^{-d_i}))$ and using
$t = q^k$ to express the central element as
$c = (1 - q)(1-t^{-1})/(q_i - q_i^{-1})$ (the standard Ding--Iohara
normalization). See~\cite{DI97} for the complete change-of-basis computation.
\end{proof}

\begin{conjecture}[Elliptic-curve toroidal realization]\label{conj:toroidal-e1}
\ClaimStatusConjectured
The toroidal algebra $U_{q,t}(\hat{\hat{\mathfrak{g}}})$ admits a
curve-based $\Eone$-chiral realization on the elliptic curve
$E_\tau$ (with modulus $\tau = \frac{\log t}{\log q}$).
\end{conjecture}

\begin{conjecture}[Double-affine surface-factorization realization]
\label{conj:toroidal-surface-factorization}
\ClaimStatusConjectured
The same toroidal data should extend to a doubly-graded
surface-factorization or double-affine object on
$\bC^* \times \bC^*$, with one direction encoding the spectral
parameter and the other the elliptic coordinate.
(This is the toroidal/elliptic extension flank of
Vol~I, Conjecture~\ref{conj:master-infinite-generator}, not the standard
$W_\infty$/Yangian finite-packet lane.)
\end{conjecture}

\begin{remark}[Scope]
Conjecture~\ref{conj:toroidal-e1} requires a direct associativity check for the $\Eone$-chiral product from the OPE relations, since the BD locality argument applies to commutative chiral algebras. The $\bC^* \times \bC^*$ realization (Conjecture~\ref{conj:toroidal-surface-factorization}) is a separate surface-factorization problem.
\end{remark}

\begin{remark}[Evidence]
We describe the conjectural $\Eone$-chiral structure.

\emph{Arrow~1: $R$-matrix braiding.}
The cross relation (2) of Definition~\ref{def:toroidal} gives the
structure function $g_{ij}(z/w) = (z/w - q^{a_{ij}})/(z/w - q^{-a_{ij}})$.
Define the $R$-matrix-valued braiding on generating currents by
\[
\sigma_{12} \circ \mu(e_i(z), e_j(w))
= g_{ij}(z/w) \cdot \mu(e_j(w), e_i(z)).
\]
Since $g_{ij}(u)$ is a rational function with $g_{ij}(u) \neq 1$
for generic $q$ and $a_{ij} \neq 0$, the braiding is non-trivial:
the algebra is associative but not commutative, hence strictly $\Eone$.

\emph{Arrow~2: Associativity (gap).}
The $R$-matrix $R_{ij}(u) = g_{ij}(u) \cdot \mathrm{Id}$ satisfies
the Yang--Baxter equation trivially (being scalar-valued).
The genuine content of associativity lies in verifying that the
full matrix-valued OPE structure, including the Serre-type relations
of Definition~\ref{def:toroidal}, defines an algebra over $\chirAss$
(Vol~I, Definition~\ref{def:chiral-ass-operad}). This requires an explicit
verification that triple OPEs are consistent, which we do not
carry out here.

\emph{Arrow~3: Realization on $E_\tau$.}
Set $\tau = \log t / \log q$ and let $E_\tau = \bC^*/q^{\bZ}$.
The generating currents $e_i(z)$, $f_i(z)$, $\psi_i(z)$ are meromorphic
on $\bC^*$ with quasi-periodicity $e_i(qz) = q^{d_i} e_i(z)$
(from the grading), so they define sections of line bundles on $E_\tau$.
The OPE of Proposition~\ref{prop:toroidal-ope} has poles at
$z/w \in q^{\bZ}$, which descend to a single point on $E_\tau$.
This defines a candidate $\Eone$-chiral algebra structure on $E_\tau$,
pending verification of Arrow~2.

\emph{Arrow~4: Surface-factorization horizon on $\bC^* \times \bC^*$.}
The horizontal modes $x^\pm_{i,n}$ depend on the spectral
parameter $u$ and the vertical modes $e_i(z)$ depend on $z$.
Together they give fields on $\bC^*_u \times \bC^*_z$ with the
two independent gradings (by horizontal and vertical mode number)
providing the doubly-graded structure. This belongs to a separate
conjectural surface-factorization or double-affine problem, not
automatically a
curve-based $\Eone$-chiral algebra. A correct formulation would
require a bona fide two-parameter factorization framework.

Conditional on Arrow~2 (i.e., the verification that the toroidal OPE
defines an associative $\chirAss$-algebra structure, which requires
checking consistency of triple OPEs with the Serre-type relations), the
$\Eone$-chiral Koszul duality framework
(Vol~I, Theorem~\ref{thm:e1-chiral-koszul-duality}) applies to the
elliptic-curve realization of Conjecture~\ref{conj:toroidal-e1}. The
$\bC^* \times \bC^*$ surface construction belongs instead to
Conjecture~\ref{conj:toroidal-surface-factorization}.
\end{remark}


\subsection{Koszul dual of the toroidal algebra}

\begin{conjecture}[Toroidal Koszul dual]
\label{conj:toroidal-koszul-dual}
\ClaimStatusConjectured
\index{toroidal algebra!Koszul dual}
The $\Eone$-chiral Koszul dual of the toroidal algebra
$U_{q,t}(\hat{\hat{\mathfrak{g}}})$ is isomorphic to
$U_{q^{-1},t^{-1}}(\hat{\hat{\mathfrak{g}}})$,
the toroidal algebra with inverted parameters.
(Contributing to the toroidal/elliptic extension flank of
Vol~I, Conjecture~\ref{conj:master-infinite-generator}.)
\end{conjecture}

\begin{remark}[Evidence]\label{rem:toroidal-koszul-evidence}
Conditional on Conjecture~\ref{conj:toroidal-e1}, the evidence is threefold:
\begin{enumerate}[label=\textup{(\roman*)}]
\item \emph{RTT presentation.} By $\Eone$-chiral Koszul duality (Vol~I, Theorem~\ref{thm:e1-chiral-koszul-duality}), the dual has $R$-matrix $R(u;q,t)^{-1} = R(u;q^{-1},t^{-1})$ (Ding--Iohara inversion; cf.\ Vol~I, Theorem~\ref{thm:yangian-koszul-dual}).
\item \emph{Affine degeneration.} In the limit $t \to 0$, $U_{q,t}(\hat{\hat{\mathfrak{g}}}) \to U_q(\hat{\mathfrak{g}})$ and the dual reduces to $U_{q^{-1}}(\hat{\mathfrak{g}})$, consistent with Vol~I, Theorem~\ref{thm:km-quantum-groups}.
\item \emph{Central charge complementarity.} The inversion $q \to q^{-1}$ sends $\tau \to -\tau$, reversing the complex structure of $E_\tau$, the elliptic analogue of the Verdier intertwining (Vol~I, Theorem~A; Theorem~\ref{thm:bar-cobar-verdier}).
\end{enumerate}
The naive elliptic Arnold relation fails (Theorem~\ref{thm:elliptic-arnold-fails}); its role passes to the dynamical Yang--Baxter coherence built on the quartic Fay identity (Theorem~\ref{thm:fay}, Conjecture~\ref{conj:dynamical-bar-nilpotency}). The universal MC class predicts $\Theta_{U_{q,t}} + \Theta_{U_{q^{-1},t^{-1}}} = 0$.
\end{remark}

\begin{remark}[Miki automorphism and Koszul duality]
\label{rem:miki-automorphism}
\index{Miki automorphism|textbf}
\index{toroidal algebra!Miki automorphism}
The quantum toroidal algebra $U_{q,t}(\hat{\hat{\fg}})$ carries an
$\SL_2(\Z)$ automorphism group, discovered by Miki~\cite{Miki07} and
implicit in the Ding--Iohara construction~\cite{DI97}. Write
$(q_1, q_2, q_3)$ for the three parameters satisfying $q_1 q_2 q_3 = 1$,
with $q = q_1$ and $t = q_2^{-1}$ in the $(q,t)$ convention. The
\emph{Miki automorphism} $S$ is the generator of $\SL_2(\Z)$ that acts
by cyclic permutation:
\[
 S \colon (q_1, q_2, q_3) \longmapsto (q_2, q_3, q_1),
\]
so that $S^3 = \mathrm{id}$ on parameters (triality). At the level of
generators, $S$ exchanges the horizontal subalgebra (generated by $E(z)$,
$F(z)$) with the vertical subalgebra (generated by $K^{\pm}(z)$):
\[
 S \colon E(z) \leftrightarrow K^+(z), \qquad
 F(z) \leftrightarrow K^-(z)
 \qquad \text{(up to normalization)}.
\]
This is the algebraic incarnation of the exchange of the two $\bC^*$
factors in the torus $\bC^* \times \bC^*$ underlying the double loop.

Three properties connect $S$ to the Koszul duality of
Conjecture~\ref{conj:toroidal-koszul-dual}:
\begin{enumerate}[label=\textup{(\alph*)}]
\item \emph{Parameter inversion.} In $(q,t)$ coordinates, $S$ acts by
 $q \mapsto 1/t$, $t \mapsto q/t$, and $S^2$ acts by
 $q \mapsto t/q$, $t \mapsto 1/q$. The composite $S^3 = \mathrm{id}$
 gives the triality. The Koszul involution $(q,t) \mapsto (q^{-1},t^{-1})$
 is separate from the six $\SL_2(\Z)$ images and commutes with the
 $\SL_2(\Z)$ action. The Koszul involution sends
 $(q_1,q_2,q_3) \mapsto (q_1^{-1},q_2^{-1},q_3^{-1})$ (total inversion,
 preserving $q_1 q_2 q_3 = 1$).

\item \emph{Affine Yangian: triality survives, full $\SL_2(\Z)$ does not.}
 The affine Yangian $Y(\hat{\fgl}_1)$ is the rational degeneration
 $q_1 \to 1$; in this limit, the three parameters collapse to
 $(\epsilon_1, \epsilon_2, \epsilon_3)$ with
 $\epsilon_1 + \epsilon_2 + \epsilon_3 = 0$, and the infinite discrete
 group $\SL_2(\Z)$ degenerates to the finite symmetric group $S_3$
 permuting $(\epsilon_1, \epsilon_2, \epsilon_3)$. This $S_3$ triality
 \emph{does} lift to algebra automorphisms of $Y(\hat{\fgl}_1)$
 (Prochazka--Rapcak); what fails in the rational limit is the survival
 of the $\Z$-extension of the Weyl group generated by the Miki element
 $S$ of infinite order. The Miki automorphism is a genuinely
 quantum-toroidal phenomenon in precisely this sense: it enhances the
rational $S_3$ to an $\SL_2(\Z)$ by placing both loops on equal
footing: $S^3 \neq \mathrm{id}$ in $U_{q,t}$ at generic $(q,t)$, so
$S$ is not a finite-order element of the rational $S_3$ triality.

\item \emph{Spectral-parameter exchange for the DDCA.} In the rational
 limit, the double deformation current algebra (DDCA) carries two
 spectral parameters $(u,v)$. The rational shadow of $S$ should exchange
 these two parameters, intertwining the horizontal and vertical RTT
 presentations. This exchange is compatible with the $R$-matrix inversion
 $R(u;q,t)^{-1} = R(u;q^{-1},t^{-1})$ appearing in item~(i) of
 Remark~\ref{rem:toroidal-koszul-evidence}: the Miki automorphism
 provides a \emph{second} route to the Koszul dual, via the internal
 symmetry of $U_{q,t}$ rather than the external bar-cobar construction.
 The precise rational degeneration identifying the DDCA as the
 $\hbar_1, \hbar_2 \to 0$ limit of the toroidal algebra is
 Conjecture~\ref{conj:ddca-toroidal-bridge}; the spectral-parameter
 exchange $\sigma \colon (u,v) \mapsto (v,u)$ of the DDCA as the
 rational shadow of $S$ is Remark~\ref{rem:ddca-miki-shadow}.
\end{enumerate}

The distinction between $S$ (internal $\SL_2(\Z)$ automorphism) and
Koszul duality (external bar-cobar inversion) is essential: $S$ is an
automorphism of $U_{q,t}$, while Koszul duality produces a
\emph{different} algebra $U_{q^{-1},t^{-1}}$. At the level of
parameters, the Koszul involution $(q_1,q_2,q_3) \mapsto
(q_1^{-1},q_2^{-1},q_3^{-1})$ commutes with the cyclic permutation
$S$, so the $\SL_2(\Z)$ symmetry of $U_{q,t}$ descends to an
$\SL_2(\Z)$ symmetry of the Koszul dual $U_{q^{-1},t^{-1}}$ with
identical parameter action.
\end{remark}

In the rational limit $q_1 \to 1$ of Remark~\ref{rem:miki-automorphism}, the Miki $\SL_2(\Z)$ collapses to the $S_3$-permutation of $(\epsilon_1, \epsilon_2, \epsilon_3)$ with $\sum \epsilon_i = 0$. This $S_3$ descends simultaneously on the Drinfeld double $Y_{\epsilon_1,\epsilon_2,\epsilon_3}(\widehat{\gl}_1)$, its positive half $Y^+ = \mathrm{CoHA}(\bC^3)$, and the Ran-factorisation coalgebra dual; the $\mathcal{W}_{1+\infty}[\lambda]$-triality of Gaiotto--Rapcak~$2019$ is the $\mathrm{ev}_\lambda$-image. The full statement with proof, together with the identification $\mathrm{CoHA}(\bC^3) = Y^+ \neq \mathcal{W}_{1+\infty}$, is Theorem~\ref{thm:plat-Miki-S3} (Vol~III, Chapter~\ref{ch:k3-quantum-toroidal}). At the Calabi--Yau slice the qq-character recursion of Nekrasov--Pestun--Shatashvili~$2018$ (arXiv:1312.6689) does not truncate: triality, not truncation, is the CY$_3$-diagonal structure (Feigin--Jimbo--Miwa--Mukhin~$2016$, arXiv:1603.02765). At the K3 promotion the $S_3$ breaks to $\Z/2\Z$ by Proposition~\ref{prop:k3-qt-no-s3-miki}: K3 admits no third algebraic-torus direction, leaving only the Fourier--Mukai involution.

\begin{conjecture}[DDCA--toroidal bridge]
\label{conj:ddca-toroidal-bridge}
\ClaimStatusConjectured
The deformed double current algebra $D(\fg)$ is the rational
degeneration of the quantum toroidal algebra
$U_{q,t}(\hat{\hat{\fg}})$: in the limit
$\hbar_1, \hbar_2 \to 0$ with $\hbar_1/\hbar_2$ fixed, the RTT
presentation of $U_{q,t}$ degenerates to the defining relations of
$D(\fg)$. Here $q=e^{\hbar_1}$ and $t=e^{-\hbar_2}$; the plain
rational parameter in the table is $\hbar=\hbar_1$ after the
$\hbar_2 \to 0$ degeneration. The intermediate degenerations are the affine Yangian
($\hbar_2 \to 0$, $\hbar_1$ fixed) and the quantum affine algebra
($\hbar_1 \to 0$, $q$ fixed).
\end{conjecture}

\begin{remark}[DDCA spectral-parameter exchange as rational Miki shadow]
\label{rem:ddca-miki-shadow}
The spectral-parameter exchange $\sigma \colon (u,v) \mapsto (v,u)$
of the DDCA is the rational shadow of the Miki automorphism~$S$
of $U_{q,t}$. Under the rational degeneration of
Conjecture~\ref{conj:ddca-toroidal-bridge}, the cyclic permutation
$S \colon (q_1, q_2, q_3) \mapsto (q_2, q_3, q_1)$ reduces to the
exchange of the two spectral parameters of $D(\fg)$.
\end{remark}

\subsection{Connection to the three main theorems of Vol~I}

\begin{remark}[Toroidal three theorems]
\label{rem:toroidal-three-theorems}
\index{three theorems!toroidal}
Conditional on Conjecture~\ref{conj:toroidal-e1}, the expected toroidal analogues of the three main theorems of Vol~I are:
\emph{Theorem~A} (Vol~I, bar-cobar adjunction): the bar complex $\bar{B}^{\mathrm{ell}}(U_{q,t})$ is an $\Eone$-chiral coassociative coalgebra once its differential carries the dynamical correction (Conjecture~\ref{conj:dynamical-bar-nilpotency}; the naive elliptic differential does not square to zero, Theorem~\ref{thm:elliptic-arnold-fails}); the Verdier chiral dual $\mathbb{D}_{\mathrm{Ran}}\bar{B}^{\mathrm{ell}}(U_{q,t}) \simeq \bar{B}^{\mathrm{ell}}(U_{q,t}^{!})$ identifies the Koszul dual algebra $U_{q,t}^{!}$.
\emph{Theorem~B} (Vol~I, completed Positselski/coderived form): the
elliptic analogue of the strict ML weight-completed counit
$\Omega\bar{B}^{\mathrm{ell}}(U_{q,t}) \xrightarrow{\sim} U_{q,t}$
is expected in the completed coderived ambient, with the rational
spectral sequence and elliptic corrections of
Theorem~\ref{thm:elliptic-vs-rational} providing the finite-stage
input; the uncompleted counit is the inversion face of Theorem~A, not
Koszul duality.
\emph{Theorem~C} (Vol~I, complementarity): $\mathrm{obs}_g(U_{q,t}) + \mathrm{obs}_g(U_{q^{-1},t^{-1}})$ controlled by $H^*$ of a framed moduli space (the $\Eone$ replacement for $\overline{\mathcal{M}}_g$; cf.\ Vol~I, Remark~\ref{rem:yangian-three-theorems}).
The universal MC class $\Theta_{U_{q,t}}$ should govern the genus expansion, with $(q,t) \mapsto (q^{-1},t^{-1})$ reflected in Verdier duality.
\end{remark}

\begin{remark}[Lattice evidence for toroidal genus theory]
\label{rem:toroidal-lattice-evidence}
\index{lattice vertex algebra!genus-1 evidence for toroidal}
The curvature-braiding orthogonality theorem for quantum lattice VOAs
(Vol~I, Theorem~\ref{thm:lattice:curvature-braiding-orthogonal}) provides
structural evidence for the toroidal case. For the lattice, the genus-$1$
Heisenberg-route curvature
$\kappa_{\mathrm{ch}}^{\mathrm{Heis}}(V_\Lambda)
= \operatorname{rank}(\Lambda)$ lives entirely in the
Cartan (Heisenberg) sector; it is not the compact CY Route~A Hodge
supertrace invariant. This Heisenberg curvature is independent of the deformation cocycle
$\varepsilon_{N,q}$. The toroidal algebra $U_{q,t}(\hat{\hat{\fg}})$
has a similar double-loop structure: the inner loop ($\hat{\fg}$) is the
Cartan analog, while the outer loop (the second affinization) provides
the braiding. The orthogonality principle predicts that the toroidal
genus-$1$ curvature should be controlled by the inner-loop data alone,
with the dynamical parameter $\lambda \in H^1(E_\tau)$ entering the
braiding axis rather than the modular axis. The dynamical YBE for
$R(u, \lambda)$ would then be the toroidal analog of the modified
Arnold relation at genus~$1$, with the dynamical shift
$\lambda \to \lambda - h^{(i)}$ replacing the B-cycle monodromy
correction.
\end{remark}

\section{Elliptic quantum groups}
\label{sec:elliptic-quantum}

The toroidal algebra acts on a double loop $\mathbb{C}^* \times \mathbb{C}^*$;
compactifying the second loop to an elliptic curve $E_\tau$ replaces the
trigonometric $R$-matrix by an elliptic one. The bare Yang--Baxter equation
$R_{12} R_{13} R_{23} = R_{23} R_{13} R_{12}$ then fails: $H^1(E_\tau) = \bC^2$
is nontrivial, so the R-matrix picks up a monodromy parameter along each
nontrivial cycle. Felder's dynamical YBE absorbs this parameter by letting
$R$ depend on a \emph{dynamical variable} $\lambda$ that shifts along the
B-cycle under propagation. The result is the elliptic quantum group
$E_{q,p}(\fg)$, in which the B-cycle monodromy of the elliptic propagator
is promoted to an algebraic parameter.

\subsection{Felder's elliptic quantum group}

\begin{definition}[Elliptic quantum group]\label{def:elliptic-quantum}
\index{elliptic quantum group|textbf}
The \emph{elliptic quantum group} $E_{q,p}(\mathfrak{g})$ (Felder) is defined by:
\begin{itemize}
\item The dynamical R-matrix $R(u, \lambda)$ depending on spectral parameter $u$
and dynamical parameter $\lambda$.
\item The RLL relations on $V \otimes V \otimes W$ ($W$ the module,
$h_W$ its Cartan action):
\[
R_{12}(u-v, \lambda - \gamma h_W)\, L_1(u, \lambda)\,
L_2(v, \lambda - \gamma h^{(1)})
= L_2(v, \lambda)\, L_1(u, \lambda - \gamma h^{(2)})\,
R_{12}(u-v, \lambda),
\]
with every dynamical shift scaled by the quantization parameter
$\gamma$, in the conventions of
Felder--Varchenko~\cite{FelderVarchenko96}.
\end{itemize}
\end{definition}

\begin{definition}[Elliptic R-matrix {\cite{Felder94,FelderVarchenko96}}]\label{def:elliptic-r}
The Felder elliptic R-matrix for $\mathfrak{sl}_2$, in the basis
$(v_+{\otimes}v_+,\, v_+{\otimes}v_-,\, v_-{\otimes}v_+,\,
v_-{\otimes}v_-)$, is
\[
R(u, \lambda) = \begin{pmatrix}
1 & 0 & 0 & 0 \\
0 & \alpha(u, \lambda) & \beta(u, \lambda) & 0 \\
0 & \beta(u, -\lambda) & \alpha(u, -\lambda) & 0 \\
0 & 0 & 0 & 1
\end{pmatrix}
\]
where $\theta = \theta_1(\cdot \mid \tau)$, $\gamma \in \bC^\times$ is the \emph{Felder quantization parameter} (disambiguated from the Arnold propagator $\eta_{ij}$ and the Weierstrass quasi-period $\eta_1$; see Remark~\ref{rem:eta-to-gamma}), and the entries are
\begin{align*}
\alpha(u, \lambda) &= \frac{\theta(u)\,\theta(\lambda + \gamma)}{\theta(u - \gamma)\,\theta(\lambda)}, \\
\beta(u, \lambda) &= -\,\frac{\theta(u + \lambda)\,\theta(\gamma)}{\theta(u - \gamma)\,\theta(\lambda)}
= -\,\frac{\theta(\gamma)}{\theta'(0)}\cdot\frac{\theta(u)}{\theta(u-\gamma)}\;
F_\tau(u, \lambda),
\end{align*}
in the common-denominator normalization $\theta(u-\gamma)\,\theta(\lambda)$ with unit corner entries. The exchange entry $\beta$ is a multiple of the Kronecker--Eisenstein kernel $F_\tau$ of Theorem~\ref{thm:fay}: the dynamical variable of the R-matrix is the second argument of the kernel.
\end{definition}

\begin{remark}[Three symbols named \texorpdfstring{$\eta$}{eta}: disambiguation]\label{rem:eta-to-gamma}
The Felder literature writes $\eta$ for the quantization parameter in $R(u,\lambda)$; the Arnold/Orlik--Solomon literature writes $\eta_{ij} = d\log\theta_1(z_i - z_j)$ for the propagator $1$-form; the Weierstrass literature writes $\eta_1 = \zeta_W(1/2;\Lambda_\tau)$ for the quasi-period. All three appear here. We reserve $\eta_{ij}$ for propagator $1$-forms and $\eta_1$ for the quasi-period, and rename the Felder parameter to $\gamma$. The translation $\eta^{\mathrm{Felder}} = \gamma$ is purely notational; every formula in Felder--Varchenko transports verbatim.
\end{remark}


\begin{remark}[Elliptic quantum groups and bar complex]
\label{rem:elliptic-qg-bar}
\index{elliptic quantum group!bar complex}
Assuming Conjecture~\ref{conj:toroidal-e1}, one expects:
(i)~$\barB(U_{q,t})$ on $E_\tau$ computes the chiral coalgebra $E_{q,p}(\mathfrak{g})$ (elliptic Kazhdan--Lusztig);
(ii)~the dynamical parameter $\lambda$ (Definition~\ref{def:elliptic-quantum}) is the B-cycle monodromy of the elliptic propagator (elliptic analogue of $k \to k + h^\vee$);
(iii)~the dynamical YBE for $R(u,\lambda)$ is the coherence law standing in for $d^2 = 0$, which fails for the naive elliptic differential (Theorem~\ref{thm:elliptic-arnold-fails}); the dynamical parameter absorbs the $B$-cycle monodromy from $H^1(E_\tau) \neq 0$ (Corollary~\ref{cor:elliptic-bar-dynamical}).
\end{remark}

\section{Elliptic R-matrices and theta functions}
\label{sec:elliptic-r-theta}

\subsection{Theta function identities}

\begin{definition}[Jacobi theta function]\label{def:theta}
The Jacobi theta function is:
\[
\theta(u | \tau) = \sum_{n \in \bZ} (-1)^n e^{\pi i \tau n(n-1) + 2\pi i n u} =
(e^{2\pi iu}; p)_\infty (p e^{-2\pi iu}; p)_\infty (p; p)_\infty
\]
where $p = e^{2\pi i \tau}$ and $(a; q)_\infty = \prod_{n \geq 0}(1 - aq^n)$.
\end{definition}

\begin{theorem}[Fay trisecant identity at genus \texorpdfstring{$1$}{1} {\cite{Fay73,WW27,BrownLevin11}}]\label{thm:fay}
\index{Fay trisecant identity|textbf}
Write $\theta = \theta_1(\cdot\,|\,\tau)$ for the odd Jacobi theta
function and
\[
F_\tau(z,\lambda)
\;=\; \frac{\theta'(0)\,\theta(z+\lambda)}{\theta(z)\,\theta(\lambda)}
\]
for the Kronecker--Eisenstein kernel. Then:
\begin{enumerate}[label=\textup{(\roman*)}]
\item \textup{(}Kernel form.\textup{)} For all $z_1, z_2$ and all
$\lambda_1, \lambda_2$ off the polar divisors,
\begin{equation}\label{eq:fay-kernel}
F_\tau(z_1,\lambda_1)\,F_\tau(z_2,\lambda_2)
\;=\;
F_\tau(z_1,\lambda_1{+}\lambda_2)\,F_\tau(z_2{-}z_1,\lambda_2)
\;+\;
F_\tau(z_2,\lambda_1{+}\lambda_2)\,F_\tau(z_1{-}z_2,\lambda_1).
\end{equation}
\item \textup{(}Quartic theta form.\textup{)} For all $a,b,c,d \in \bC$,
\begin{equation}\label{eq:fay-quartic}
\theta(a{+}b)\,\theta(a{-}b)\,\theta(c{+}d)\,\theta(c{-}d)
+ \theta(b{+}c)\,\theta(b{-}c)\,\theta(a{+}d)\,\theta(a{-}d)
+ \theta(c{+}a)\,\theta(c{-}a)\,\theta(b{+}d)\,\theta(b{-}d)
\;=\; 0,
\end{equation}
and the same identity holds verbatim for the Weierstrass sigma
function of $\Lambda_\tau = \bZ + \tau\bZ$.
\end{enumerate}
Every term of either form is a product of \emph{four} theta factors;
the quartic structure is forced. The bilinear expression
\begin{equation}\label{eq:bilinear-three-term}
\theta(u_1{-}u_2)\,\theta(u_3{-}u_4)
- \theta(u_1{-}u_3)\,\theta(u_2{-}u_4)
+ \theta(u_1{-}u_4)\,\theta(u_2{-}u_3)
\end{equation}
is not identically zero for any $\tau$
(Lemma~\ref{lem:no-bilinear-fay}); it vanishes only in the
trigonometric degeneration $\tau \to i\infty$, in which (after
removing the overall factor $q^{1/4}$) $\theta_1(z|\tau) \to
2\sin(\pi z)$ and \eqref{eq:bilinear-three-term} becomes the Ptolemy
identity $\sin(\alpha{-}\beta)\sin(\gamma{-}\delta) -
\sin(\alpha{-}\gamma)\sin(\beta{-}\delta) +
\sin(\alpha{-}\delta)\sin(\beta{-}\gamma) = 0$; the further limit
$\sin(\pi z) \to \pi z$ recovers the rational partial-fraction
identity.
\end{theorem}

\begin{proof}
\emph{Kernel form.} Fix $z_2, \lambda_1, \lambda_2$ generic and let
$G(z_1)$ be the difference of the two sides
of~\eqref{eq:fay-kernel}. From $\theta(z+1) = -\theta(z)$ and
$\theta(z+\tau) = -e^{-\pi i\tau - 2\pi i z}\theta(z)$ the kernel
transforms by $F_\tau(z+1,\lambda) = F_\tau(z,\lambda)$ and
$F_\tau(z+\tau,\lambda) = e^{-2\pi i\lambda} F_\tau(z,\lambda)$.
Each of the three products in~\eqref{eq:fay-kernel} acquires the same
factor $e^{-2\pi i\lambda_1}$ under $z_1 \mapsto z_1 + \tau$ (for the
middle term, $e^{-2\pi i(\lambda_1+\lambda_2)}$ from its first factor
times $e^{+2\pi i\lambda_2}$ from its second) and is invariant under
$z_1 \mapsto z_1 + 1$. The candidate poles of $G$ cancel: at
$z_1 = 0$ the first and second products both have residue
$F_\tau(z_2,\lambda_2)$ and enter $G$ with opposite signs; at
$z_1 = z_2$ the second and third products have residues
$-F_\tau(z_2,\lambda_1{+}\lambda_2)$ and
$+F_\tau(z_2,\lambda_1{+}\lambda_2)$ (the orientation reversal in
$F_\tau(z_2{-}z_1,\lambda_2)$ flips the sign) and enter $G$ with the
same sign. So $G$ is a holomorphic section of a degree-zero line
bundle on $E_\tau$ with multiplier $e^{-2\pi i\lambda_1} \neq 1$ for
$\lambda_1 \notin \Lambda_\tau$; every such section vanishes
identically. Analytic continuation removes the genericity.

\emph{Quartic form.} Clearing denominators in~\eqref{eq:fay-kernel}
(multiply by $\theta(z_1)\theta(z_2)\theta(\lambda_1)
\theta(\lambda_2)\theta(\lambda_1{+}\lambda_2)\theta(z_1{-}z_2)
/\theta'(0)^2$ and use oddness of $\theta$) gives
\begin{multline}\label{eq:fay-quartic-cleared}
\theta(z_1{+}\lambda_1)\,\theta(z_2{+}\lambda_2)\,
\theta(\lambda_1{+}\lambda_2)\,\theta(z_1{-}z_2)
+ \theta(z_1{+}\lambda_1{+}\lambda_2)\,
\theta(\lambda_2{+}z_2{-}z_1)\,\theta(\lambda_1)\,\theta(z_2) \\
- \theta(z_2{+}\lambda_1{+}\lambda_2)\,
\theta(\lambda_1{+}z_1{-}z_2)\,\theta(\lambda_2)\,\theta(z_1) \;=\; 0,
\end{multline}
which is~\eqref{eq:fay-quartic} under the invertible affine
substitution
$a = \tfrac{1}{2}(z_1{+}\lambda_1{+}z_2{+}\lambda_2)$,
$b = \tfrac{1}{2}(z_1{+}\lambda_1{-}z_2{-}\lambda_2)$,
$c = \tfrac{1}{2}(\lambda_1{+}\lambda_2{+}z_1{-}z_2)$,
$d = \tfrac{1}{2}(\lambda_1{+}\lambda_2{-}z_1{+}z_2)$
(oddness of $\theta$ pairs up the signs). The sigma form follows from
$\sigma(z) = e^{\eta_1 z^2}\,\theta(z)/\theta'(0)$: in each term
of~\eqref{eq:fay-quartic} the exponential factors multiply to the
common factor $e^{2\eta_1(a^2+b^2+c^2+d^2)}$.
\end{proof}

\begin{lemma}[No bilinear trisecant at genus \texorpdfstring{$1$}{1}]
\label{lem:no-bilinear-fay}
For every $\tau \in \mathbb{H}$ the
expression~\eqref{eq:bilinear-three-term} is not identically zero in
$(u_1, u_2, u_3, u_4)$.
\end{lemma}

\begin{proof}
Suppose it vanished identically. Fix generic $u_2, u_3, u_4$ and read
the vanishing as a linear relation among the three functions
$f_j(u_1) = \theta(u_1 - u_j)$, $j = 2,3,4$, with the nonzero
coefficients $\theta(u_3{-}u_4)$, $-\theta(u_2{-}u_4)$,
$\theta(u_2{-}u_3)$. Applying $u_1 \mapsto u_1 + \tau$ once and twice
and using $\theta(z+\tau) = -e^{-\pi i\tau - 2\pi i z}\theta(z)$
produces two further relations whose coefficients are rescaled by
$e^{2\pi i u_j}$ and $e^{4\pi i u_j}$. The $3 \times 3$ coefficient
matrix of the three relations is Vandermonde in the distinct values
$e^{2\pi i u_2}, e^{2\pi i u_3}, e^{2\pi i u_4}$, hence invertible;
so each $f_j$ vanishes identically --- false. In the trigonometric
degeneration the three multipliers collapse to one and the
obstruction disappears; this is why the Ptolemy identity has no
elliptic lift.
\end{proof}


\subsection{Failure of the elliptic Arnold relation}

What happens to the Arnold relation when the propagator becomes
doubly periodic? The rational relation rests on the partial-fraction
identity $(z_1{-}z_2)^{-1}(z_2{-}z_3)^{-1} + \text{cyclic} = 0$; its
elliptic candidate replaces $(z_i - z_j)^{-1}$ by
$\zeta_\tau(z_i - z_j) := \partial_z \log\theta_1(z_i - z_j|\tau)$
(the logarithmic derivative of $\theta_1$; this differs from the
classical Weierstrass $\zeta$-function $\zeta_W(z) = \zeta_\tau(z)
+ 2\eta_1 z$ by a linear term, where $\eta_1 = \zeta_W(1/2)$). The
candidate fails, and the exact size of the failure is a theorem. The
classical identity that does hold localizes the \emph{square of the
sum}, not the sum of pairwise products.

\begin{proposition}[Weierstrass square identity]\label{prop:zeta-square}
\index{Weierstrass zeta function!square identity}
Let $\zeta_W$ and $\wp$ be the Weierstrass functions of
$\Lambda_\tau$ and let $a + b + c = 0$ with $a, b, c \notin
\Lambda_\tau$. Then
\begin{equation}\label{eq:zeta-square}
\bigl(\zeta_W(a) + \zeta_W(b) + \zeta_W(c)\bigr)^2
\;=\; \wp(a) + \wp(b) + \wp(c).
\end{equation}
The left side is unchanged by the substitution
$\zeta_W \to \zeta_\tau$, since $\zeta_\tau(z) = \zeta_W(z) -
2\eta_1 z$ and the linear terms cancel on $a + b + c = 0$.
\end{proposition}

\begin{proof}
The addition theorems
$\zeta_W(a{+}b) = \zeta_W(a) + \zeta_W(b) +
\tfrac{1}{2}\,\frac{\wp'(a)-\wp'(b)}{\wp(a)-\wp(b)}$
and
$\wp(a{+}b) = \tfrac{1}{4}\bigl(\frac{\wp'(a)-\wp'(b)}
{\wp(a)-\wp(b)}\bigr)^2 - \wp(a) - \wp(b)$
combine: with $c = -(a{+}b)$, oddness of $\zeta_W$ gives
$\sum\zeta_W = -\tfrac{1}{2}(\wp'(a)-\wp'(b))/(\wp(a)-\wp(b))$;
squaring and substituting the second theorem, with
$\wp(c) = \wp(a{+}b)$ by evenness, yields~\eqref{eq:zeta-square}.
\end{proof}

\begin{theorem}[Failure of the naive elliptic Arnold relation]
\label{thm:elliptic-arnold-fails}
\index{Arnold relation!failure at genus 1|textbf}
On $C_3(E_\tau)$ write $a = z_1 - z_2$, $b = z_2 - z_3$,
$c = z_3 - z_1$ (so $a + b + c = 0$) and let
$\eta_{ij}^{\mathrm{ell}} = \zeta_\tau(z_i - z_j)\,(dz_i - dz_j)$.
Then:
\begin{enumerate}[label=\textup{(\roman*)}]
\item The cyclic wedge sum is a multiple of $da \wedge db$ whose
scalar coefficient is the cyclic sum $S_\tau$:
\[
\eta_{12}^{\mathrm{ell}} \wedge \eta_{23}^{\mathrm{ell}}
+ \eta_{23}^{\mathrm{ell}} \wedge \eta_{31}^{\mathrm{ell}}
+ \eta_{31}^{\mathrm{ell}} \wedge \eta_{12}^{\mathrm{ell}}
= S_\tau(a,b)\; da \wedge db,
\quad
S_\tau = \zeta_\tau(a)\zeta_\tau(b) + \zeta_\tau(b)\zeta_\tau(c)
+ \zeta_\tau(c)\zeta_\tau(a).
\]
\item \textup{(}Exact defect formula.\textup{)}
$\displaystyle
2\,S_\tau
= \Bigl(\sum \zeta_\tau\Bigr)^{2} - \sum \zeta_\tau^{2}
= \sum_{x \in \{a,b,c\}} \bigl(\wp(x) - \zeta_\tau(x)^2\bigr)$,
by Proposition~\ref{prop:zeta-square}.
\item For every $\tau \in \mathbb{H}$, $S_\tau$ is a nonconstant
meromorphic function on $\{a+b+c=0\}$; in particular
$S_\tau \not\equiv 0$: the naive elliptic Arnold relation
\emph{fails}. Its collision limit is quasi-modular:
\[
\lim_{(a,b)\to 0} S_\tau \;=\; 6\,\eta_1 \;=\; \pi^2 E_2(\tau).
\]
\item In the degenerations: on the cylinder
($\zeta_{\mathrm{trig}}(z) = \pi\cot(\pi z)$) the defect is the
\emph{constant} $S_{\mathrm{trig}} \equiv \pi^2$; on the line
($\zeta_{\mathrm{rat}}(z) = 1/z$) the defect vanishes,
$S_{\mathrm{rat}} \equiv 0$, and the Arnold relation
(Vol~I, Theorem~\ref{thm:arnold-relations}) is restored.
\end{enumerate}
\end{theorem}

\begin{proof}
(i) With $da = dz_1 - dz_2$, $db = dz_2 - dz_3$, $dc = -da - db$ one
has $db \wedge dc = da \wedge db = dc \wedge da$, so the three wedge
terms share the factor $da \wedge db$ with the stated scalar
coefficients.

(ii) $2S_\tau = (\sum\zeta_\tau)^2 - \sum\zeta_\tau^2$; the first
summand equals $(\sum\zeta_W)^2 = \sum\wp$ by
Proposition~\ref{prop:zeta-square}.

(iii) Let $a \to 0$ at fixed generic $b$ (so $c \to -b$). Then
$\zeta_\tau(b) + \zeta_\tau(c) = \zeta_\tau(b) - \zeta_\tau(b{+}a)
= -a\,\zeta_\tau'(b) + O(a^2)$, so
$\zeta_\tau(a)\bigl(\zeta_\tau(b)+\zeta_\tau(c)\bigr) \to
-\zeta_\tau'(b) = \wp(b) + 2\eta_1$, while
$\zeta_\tau(b)\zeta_\tau(c) \to -\zeta_\tau(b)^2$; hence
\[
\lim_{a \to 0} S_\tau \;=\; \wp(b) + 2\eta_1 - \zeta_\tau(b)^2 .
\]
If $S_\tau$ were constant, this limit would be constant in $b$,
making $\zeta_\tau^2$ elliptic; but the $B$-cycle monodromy of
$d\log\theta_1$ gives $\zeta_\tau(z + \tau) = \zeta_\tau(z) - 2\pi i$,
so $\zeta_\tau(z{+}\tau)^2 - \zeta_\tau(z)^2 = -4\pi i\,\zeta_\tau(z)
- 4\pi^2$ is nonconstant while $\wp$ is elliptic --- a contradiction,
for every $\tau$. For the collision limit, insert the Laurent
expansion $\zeta_\tau(x) = 1/x - 2\eta_1 x + O(x^3)$ (Arrow~1 of the
proof of Theorem~\ref{thm:elliptic-vs-rational}): on $a + b + c = 0$,
\[
S_\tau = \Bigl(\tfrac{1}{ab} + \tfrac{1}{bc} + \tfrac{1}{ca}\Bigr)
- 2\eta_1 \sum_{x \neq y}\tfrac{x}{y} + O(|a|^2+|b|^2)
= 0 + 6\eta_1 + O(|a|^2+|b|^2),
\]
using $\tfrac{1}{ab} + \tfrac{1}{bc} + \tfrac{1}{ca} =
\tfrac{a+b+c}{abc} = 0$ and $\sum_{x \neq y} x/y =
-(a^3{+}b^3{+}c^3)/(abc) = -3$ (since $a^3+b^3+c^3 = 3abc$ on
$a+b+c=0$), together with $\eta_1 = \tfrac{\pi^2}{6}E_2(\tau)$.

(iv) Trigonometric: for $A+B+C = 0$,
$\cot B\cot C + \cot C\cot A + \cot A\cot B = 1$ (from
$\cot C = (1 - \cot A\cot B)/(\cot A + \cot B)$), so
$S_{\mathrm{trig}} = \pi^2$. Rational: partial fractions, as in
Vol~I, Theorem~\ref{thm:bar-nilpotency-complete}.
\end{proof}

\begin{corollary}[The elliptic bar differential requires the
dynamical correction]\label{cor:elliptic-bar-dynamical}
\index{bar complex!elliptic!dynamical correction required}
Let $d^{\mathrm{ell}}$ be the elliptic bar differential built from
the propagator $\eta^{\mathrm{ell}}_{ij}$ alone, following the
rational construction of
Vol~I, \S\ref{sec:bar-nilpotency-nine-terms-complete}. The
double-contraction terms of $(d^{\mathrm{ell}})^2$ on
$\overline{C}_3(E_\tau)$ carry the coefficient
$S_\tau\, da \wedge db \neq 0$ of
Theorem~\ref{thm:elliptic-arnold-fails}, so the rational cancellation
mechanism is unavailable: the naive elliptic bar differential does
not square to zero. Two repairs are available.
\begin{enumerate}[label=\textup{(\roman*)}]
\item \emph{Central absorption (abelian targets).} When the double
contractions land in the centre --- the Heisenberg class --- the
defect is realised as the curvature $m_0^{\mathrm{ell}} \propto
E_2(\tau)$ of Computation~\ref{comp:ell-curvature}, the quasi-modular
collision value of $S_\tau$: the bar object is a \emph{curved}
coalgebra, and $d^2 = 0$ holds only after the curved twist. The
trigonometric constant $S_{\mathrm{trig}} = \pi^2$ is the
$\tau \to i\infty$ shadow.
\item \emph{Dynamical correction (general targets).} A
point-dependent defect cannot be absorbed by a central term. The
differential must be corrected by Felder's dynamical parameter
$\lambda$ (\S\ref{sec:elliptic-quantum}): the $B$-cycle monodromy
$\zeta_\tau(z+\tau) = \zeta_\tau(z) - 2\pi i$ responsible for
$S_\tau \neq 0$ is the monodromy that the dynamical shift
$\lambda \mapsto \lambda - \gamma h^{(i)}$ tracks, and the coherence
law standing in for $d^2 = 0$ is the dynamical Yang--Baxter equation
(Definition~\ref{def:dybe},
Conjecture~\ref{conj:dynamical-bar-nilpotency}). The identity that
survives at genus~$1$ is the kernel-form Fay
identity~\eqref{eq:fay-kernel}, in which the dynamical variable is
already present; a $\lambda$-free bilinear identity does not exist
(Lemma~\ref{lem:no-bilinear-fay}).
\end{enumerate}
\end{corollary}

\begin{remark}[Degeneration of the defect]\label{rem:fay-to-arnold}
Under the two-step degeneration of
Remark~\ref{rem:degeneration-limits} the defect of
Theorem~\ref{thm:elliptic-arnold-fails} descends the ladder
\[
S_\tau(a,b) \;\xrightarrow{\;\tau \to i\infty\;}\;
S_{\mathrm{trig}} = \pi^2 \;\xrightarrow{\;q \to 1\;}\;
S_{\mathrm{rat}} = 0 :
\]
point-dependent and quasi-modular at genus~$1$ (dynamical parameter
required), constant on the cylinder (absorbable as central
curvature), zero on the line (Arnold restored, and with it the bar
nilpotency of Vol~I, Theorem~\ref{thm:bar-nilpotency-complete}).
The collision value $\pi^2 E_2(\tau) \to \pi^2$ as $E_2 \to 1$ makes
the ladder consistent; under the rational rescaling
$z \mapsto \varepsilon z$ the constant cylinder defect is subleading
against the $1/\varepsilon^2$ propagator singularity and drops out
exactly.
\end{remark}


\subsection{Bar complex with theta functions}

\begin{construction}[Elliptic bar complex]\label{constr:elliptic-bar}
The bar complex of an elliptic chiral algebra incorporates theta functions:

\emph{Forms.} Replace $d\log(z_1 - z_2)$ with:
\[
\eta_{12} = d\log\theta(z_1 - z_2 | \tau)
\]
the logarithmic derivative of theta.

\emph{Differential.} The residue is computed at the zeros of $\theta$:
\[
d_{\text{res}}([a|b] \otimes \eta_{12}) = \sum_{\theta(z_0) = 0} \Res_{z_1 = z_0 + z_2}[\cdots]
\]
The zeros of $\theta(u|\tau)$ are at $u = m + n\tau$ for $m, n \in \bZ$ (the lattice points).
\end{construction}

\begin{theorem}[Elliptic vs rational homology]\label{thm:elliptic-vs-rational}
Let $\cA$ be a rational $C_2$-cofinite vertex/chiral algebra whose
genus-$1$ correlation functions satisfy Zhu modularity, and assume
the nilpotency package \textup{(H$_{\mathrm{nil}}$)}: the elliptic
bar differential on $\B^{\mathrm{ell}}(\cA)$ squares to zero. The
package is not automatic --- the naive differential fails to square
to zero (Theorem~\ref{thm:elliptic-arnold-fails}); it is supplied
for abelian (Heisenberg-class) $\cA$ by the curved twist of
Corollary~\ref{cor:elliptic-bar-dynamical}(i), and for general
targets it is the dynamical extension of
Conjecture~\ref{conj:dynamical-bar-nilpotency}, whose construction
is the named obligation. Under \textup{(H$_{\mathrm{nil}}$)}, the
homology of the elliptic bar complex on $E_\tau$ carries a finite
correction-order filtration
\[
0 = \mathcal F^{n+1}H_n \subset \mathcal F^nH_n \subset \cdots
\subset \mathcal F^1H_n \subset \mathcal F^0H_n
= H_n(\B^{\mathrm{ell}}(\cA)),
\]
with
\[
\operatorname{gr}^0_{\mathcal F}H_n \cong H_n(\B^{\mathrm{rat}}(\cA)).
\]
For $k\geq 1$, the correction quotient
$\operatorname{gr}^k_{\mathcal F}H_n$ is a finite $\cA$-module whose
$\tau$-dependence lies in quasi-modular forms of total weight~$2k$.
Specifically:
\begin{enumerate}[label=(\roman*)]
\item The leading correction is controlled by $E_2(\tau)$, the
weight-$2$ quasi-modular Eisenstein series, arising from the
regularisation of
the elliptic propagator $\partial_z \log\theta(z|\tau)$
relative to the rational propagator $1/z$
(cf.\ Construction~\ref{constr:elliptic-bar}).
\item Higher correction quotients are generated over the
quasi-modular ring by products of the coefficients
$c_j(\tau)\in M_{2j}(\mathrm{SL}_2(\bZ))$ for $j\geq 2$ and
$c_1(\tau)\in\widetilde M_2(\mathrm{SL}_2(\bZ))$; the first
holomorphic modular coefficient is the $E_4(\tau)$ term.
\item The conformal dimension $h$ of the states entering the
bar complex constrains the weights: a state of dimension $h$
contributes to $\operatorname{gr}^k_{\mathcal F}H_n$ only for
$k \leq h+n$.
\item At the rational limit $\tau \to i\infty$ (i.e., $q = e^{2\pi i\tau} \to 0$),
the Eisenstein series $E_{2k}(\tau) \to 1$ and all corrections
degenerate to constants, recovering
$H_n(\B^{\mathrm{rat}}(\cA))$ as the associated graded leading term,
up to scalar renormalisation.
\end{enumerate}
For the Heisenberg algebra ($c = 1$), the corrections reduce to the
Eisenstein series hierarchy computed in
Vol~I, \S\ref{sec:heisenberg-genus-expansion-eisenstein}.
\end{theorem}

\begin{proof}
We construct the filtration via a perturbative spectral sequence
comparing the elliptic and rational bar differentials.

\emph{Arrow~1 (Propagator decomposition).}
The elliptic bar complex (Construction~\ref{constr:elliptic-bar})
uses the propagator $\eta_{ij}^{\mathrm{ell}} = d\log\theta_1(z_i -
z_j|\tau)$, while the rational bar complex uses $\eta_{ij}^{\mathrm{rat}}
= d\log(z_i - z_j)$. Near $z_i = z_j$, both propagators have the same
simple pole with residue~$1$, so their difference $\delta_{ij}(\tau) :=
\eta_{ij}^{\mathrm{ell}} - \eta_{ij}^{\mathrm{rat}}$ is a regular
$1$-form. The logarithmic derivative $\zeta_\tau(z) :=
\partial_z\log\theta_1(z|\tau)$ differs from the classical
Weierstrass $\zeta$-function $\zeta_W(z;\Lambda_\tau)$ by
$\zeta_\tau(z) = \zeta_W(z) - 2\eta_1 z$ where
$\eta_1 = \zeta_W(1/2;\Lambda_\tau)$. Using the
Eisenstein--Weierstrass expansion of $\zeta_W$ and subtracting
the linear term gives the Laurent expansion of $\zeta_\tau$:

\[
\partial_z\log\theta_1(z|\tau)
\;=\; \frac{1}{z} + \sum_{k=1}^{\infty} c_k(\tau)\,z^{2k-1},
\]
where the correction coefficients $c_k(\tau)$ are expressed in terms
of Eisenstein series via $G_{2k}(\tau) = 2\zeta_{\mathrm{R}}(2k)\,
E_{2k}(\tau)$:
\begin{itemize}
\item $c_1(\tau) = -\tfrac{\pi^2}{3}E_2(\tau) \in
\widetilde{M}_2(\mathrm{SL}_2(\bZ))$ (quasi-modular, weight~$2$),
since $c_1 = -2\eta_1$ where
$\eta_1 = \zeta(1/2;\Lambda_\tau) = \tfrac{\pi^2}{6}E_2(\tau)$;
\item $c_k(\tau) = -G_{2k}(\tau) \in M_{2k}(\mathrm{SL}_2(\bZ))$
for $k \geq 2$ (holomorphic modular forms).
\end{itemize}
The expansion converges absolutely for $|z| < \min_{\omega \in
\Lambda_\tau \setminus\{0\}} |\omega|$.

\emph{Arrow~2 (Correction-order filtration).}
Define the \emph{correction-order filtration}
$\mathcal{F}^p \bar{B}^{\mathrm{ell}}_n(\cA)$ by
declaring that an element has filtration degree $\geq p$ if it arises
from at least $p$ insertions of the regular correction terms
$\{c_k(\tau)\,z^{2k-1}\}_{k \geq 1}$ in the propagator expansion.
Equivalently, replace each propagator by $1/z + t\cdot(c_1 z + c_2 z^3
+ \cdots)$ and expand in powers of the formal parameter~$t$; then
$\mathcal{F}^p$ consists of terms of order~$\geq t^p$.

Since the bar differential $d^{\mathrm{ell}}$ at degree~$n$
involves at most $n$ propagator factors, we have $\mathcal{F}^{n+1}
\bar{B}^{\mathrm{ell}}_n = 0$. The filtration is compatible
with the differential: $d^{\mathrm{ell}}(\mathcal{F}^p) \subset
\mathcal{F}^p$ (inserting the bar differential does not decrease
the number of correction insertions).

\emph{Arrow~3 (Spectral sequence: $E_0$ and $E_1$ pages).}
The associated graded $\mathrm{gr}^0_{\mathcal{F}}
\bar{B}^{\mathrm{ell}}(\cA)$ is the bar complex computed using only the
singular part $1/z$ of the propagator. Since the residue extraction
in the bar differential (which encodes the OPE data of~$\cA$) depends
only on the polar part, the differential on $\mathrm{gr}^0$ coincides
with the rational bar differential $d^{\mathrm{rat}}$. Thus:
\[
E_1^{0,n} = H_n(\bar{B}^{\mathrm{rat}}(\cA)).
\]
More generally, $E_1^{p,n} = H_n(\bar{B}^{\mathrm{rat}}(\cA)) \otimes
\mathrm{Sym}^p_{\geq 1}\{c_k(\tau)\}$, where the symmetric product is
taken over all propagator correction coefficients contributing total
weight~$p$.

\emph{Arrow~4 (Modular structure via Zhu's theorem).}
The differential $d_r \colon E_r^{p,n} \to E_r^{p+r,n+1}$ on the
$r$-th page arises from the interaction of $r$ correction insertions
with the bar differential. By Zhu's modular invariance
theorem~\cite{Zhu96}, the $n$-point genus-$1$ correlation functions
$\langle V_1(z_1) \cdots V_n(z_n) \rangle_{E_\tau}$ of a rational
vertex algebra $\cA$ transform as components of a vector-valued modular
form for $\mathrm{SL}_2(\bZ)$. Since the bar complex at degree~$n$ is
built from such correlation functions composed with the propagator form,
and the correction coefficients $c_k(\tau)$ are themselves
(quasi-)modular, the differentials $d_r$ respect the modular weight
grading. Consequently, the $E_\infty$ terms inherit a weight
decomposition from the ring $\widetilde{M}_*(\mathrm{SL}_2(\bZ))$
of quasi-modular forms.

\emph{Arrow~5 (Convergence).}
The spectral sequence converges because the filtration is bounded:
$\mathcal{F}^{n+1}\bar{B}^{\mathrm{ell}}_n = 0$ at each bar
degree~$n$. At convergence, $H_n(\bar{B}^{\mathrm{ell}}(\cA))$
admits the filtration
\[
0 = \mathcal{F}^{n+1} H_n \subset \mathcal{F}^n H_n \subset \cdots
\subset \mathcal{F}^1 H_n \subset \mathcal{F}^0 H_n
= H_n(\bar{B}^{\mathrm{ell}}(\cA))
\]
with successive quotients $\mathcal{E}^{(k)}_n := E_\infty^{k,n}$.
The modular weight grading controls the associated graded pieces. A
choice of weight-preserving splitting over~$\bC$ gives a vector-space
decomposition, but the canonical statement is the filtration and its
associated graded.

Properties~(i)--(iv) now follow from the structure of the correction
coefficients:
\begin{enumerate}[label=(\roman*)]
\item The leading correction involves $c_1(\tau) = -(\pi^2/3)\,E_2(\tau)$
(quasi-modular, weight~$2$).
\item Higher corrections involve products of $c_1(\tau)$ and of
$c_k(\tau) \in M_{2k}(\mathrm{SL}_2(\bZ))$ for $k \geq 2$.
\item A bar element of degree~$n$ involves at most $n$ propagator
factors; a state of conformal dimension~$h$ constrains the Laurent
order of each propagator insertion to $\leq 2h - 1$, giving
$\operatorname{gr}^k_{\mathcal F}H_n=0$ for $k > h+n$.
\item As $\tau \to i\infty$: $q \to 0$ and $E_{2k}(\tau) \to 1$, so all
corrections reduce to constants, recovering $H_n(\bar{B}^{\mathrm{rat}}(\cA))$
up to renormalization.
\end{enumerate}

\emph{Verification for the Heisenberg algebra.}
For $\cA = \mathcal{H}$ ($c = 1$), the genus-$1$ two-point function is
$\langle a(z_1)a(z_2)\rangle_{E_\tau} = \kappa_{\mathrm{ch}}\,G_\tau(z_1 - z_2)$
with $G_\tau(z) = \zeta_\tau(z) - (\pi^2 E_2(\tau)/3)\,z$
(Vol~I, Theorem~\ref{thm:heisenberg-genus-one-complete}). The leading
correction quotient $\operatorname{gr}^1_{\mathcal F}H_n$ arises from the
$c_1$ term. The complete Eisenstein hierarchy of
Vol~I, \S\ref{sec:heisenberg-genus-expansion-eisenstein} gives
the higher coefficients for descendants and higher-pole insertions.
\end{proof}

\begin{remark}[Scope]
The proof assembles three ingredients:
(i)~the Laurent expansion of the Weierstrass $\zeta$-function in
terms of Eisenstein series (classical);
(ii)~the convergence of the correction-order spectral sequence
(bounded filtration at each bar degree); and
(iii)~Zhu's modular invariance theorem for genus-$1$ $n$-point
functions of rational vertex algebras~\cite{Zhu96}. The output is the
correction-order filtration above; a direct-sum decomposition requires
a non-canonical choice of splitting.
The Heisenberg verification uses the explicit computations of
Vol~I, \S\ref{sec:heisenberg-genus-expansion-eisenstein}.
The nilpotency hypothesis \textup{(H$_{\mathrm{nil}}$)} is a genuine
restriction: without the curved twist or the dynamical correction
there is no complex whose homology the statement could describe
(Theorem~\ref{thm:elliptic-arnold-fails}).
\end{remark}



\section{Explicit elliptic bar complex for the Heisenberg algebra}
\label{sec:elliptic-bar-heisenberg}
\index{bar complex!elliptic Heisenberg|textbf}
\index{Heisenberg algebra!elliptic bar complex}

\subsection{Bar elements and the elliptic propagator}
\label{subsec:elliptic-bar-elements}

We work on $E_\tau = \bC / (\bZ + \tau\bZ)$ with the elliptic
propagator
\begin{equation}\label{eq:ell-propagator}
\eta_{ij}^{\mathrm{ell}}
= d\log\theta_1(z_i - z_j \,|\, \tau)
= \zeta_\tau(z_i - z_j)\,(dz_i - dz_j),
\end{equation}
where $\zeta_\tau(z) = \partial_z \log\theta_1(z|\tau)$ is the
logarithmic derivative of $\theta_1$ on $E_\tau$
(see the discussion preceding Proposition~\ref{prop:zeta-square}
for the distinction from the classical Weierstrass
$\zeta$-function).

\begin{computation}[Degree-1 bar elements]\label{comp:ell-bar-deg1}
\index{bar complex!elliptic!degree 1}
In bar degree~$1$, the generators of $\B^{\mathrm{ell}}_1(\cH_{\kappa_{\mathrm{ch}}})$
are:
\[
[a_m] \otimes 1 \in \B^{\mathrm{ell}}_1(\cH_{\kappa_{\mathrm{ch}}}), \quad m \in \bZ,
\]
where $a_m$ denotes the mode of the Heisenberg current $a(z) = \sum_m a_m z^{-m-1}$.
The bar differential on degree-$1$ elements is:
\[
d^{\mathrm{ell}}([a_m] \otimes 1) = 0,
\]
since the Heisenberg algebra has trivial unary OPE
($a_{(n)}a = 0$ for $n \geq 1$ except $a_{(1)}a = \kappa_{\mathrm{ch}}$, and
the residue extraction on a single element yields zero).
This is identical to the rational case.
\end{computation}

\begin{computation}[Degree-2 bar elements and the elliptic differential]
\label{comp:ell-bar-deg2}
\index{bar complex!elliptic!degree 2}
In bar degree~$2$, the generators are:
\[
[a_m | a_n] \otimes \eta_{12}^{\mathrm{ell}}
\in \B^{\mathrm{ell}}_2(\cH_{\kappa_{\mathrm{ch}}}),
\]
where $\eta_{12}^{\mathrm{ell}} = \zeta_\tau(z_1 - z_2)\,(dz_1 - dz_2)$.
The bar differential acts by residue extraction:
\begin{align}\label{eq:ell-bar-diff-deg2}
d^{\mathrm{ell}}([a_m | a_n] \otimes \eta_{12}^{\mathrm{ell}})
&= \Res_{z_1 = z_2}\Bigl[
a_m(z_1)\,a_n(z_2)\,\eta_{12}^{\mathrm{ell}}
\Bigr] \notag \\
&= \Res_{\epsilon = 0}\left[
\frac{\kappa_{\mathrm{ch}}}{\epsilon^2}\,
\zeta_\tau(\epsilon)\,d\epsilon
\right] \cdot z_2^{-m-n-2}\,dz_2,
\end{align}
where $\epsilon = z_1 - z_2$. We use the Laurent expansion of
$\zeta_\tau$ (established in the proof of
Theorem~\ref{thm:elliptic-vs-rational}, Arrow~1):
\begin{equation}\label{eq:zeta-laurent}
\zeta_\tau(\epsilon)
= \frac{1}{\epsilon}
- \frac{\pi^2}{3}\,E_2(\tau)\,\epsilon
- \frac{\pi^4}{45}\,E_4(\tau)\,\epsilon^3
- \frac{2\pi^6}{945}\,E_6(\tau)\,\epsilon^5
- \cdots
\end{equation}
Substituting into the residue:
\begin{align}\label{eq:ell-residue-expansion}
\Res_{\epsilon = 0}\left[
\frac{\kappa_{\mathrm{ch}}}{\epsilon^2}\,
\zeta_\tau(\epsilon)\,d\epsilon
\right]
&= \Res_{\epsilon = 0}\left[
\frac{\kappa_{\mathrm{ch}}}{\epsilon^2}\left(
\frac{1}{\epsilon}
- \frac{\pi^2}{3}\,E_2(\tau)\,\epsilon
- \cdots
\right) d\epsilon
\right] \notag \\
&= \Res_{\epsilon = 0}\left[
\frac{\kappa_{\mathrm{ch}}}{\epsilon^3}
- \frac{\kappa_{\mathrm{ch}}\pi^2}{3}\,E_2(\tau)\,\frac{1}{\epsilon}
- \cdots
\right] d\epsilon \notag \\
&= -\frac{\kappa_{\mathrm{ch}}\pi^2}{3}\,E_2(\tau).
\end{align}
The $1/\epsilon^3$ term has zero residue; the first nonzero
contribution is the $E_2(\tau)$ term. In the rational case
($\zeta_\tau(\epsilon) \to 1/\epsilon$), the residue of
$\kappa_{\mathrm{ch}}\epsilon^{-3}\,d\epsilon$ vanishes, so the rational
bar differential on this element is zero. The
\emph{entire degree-$2$ bar differential is the elliptic
correction}:
\begin{equation}\label{eq:ell-bar-diff-result}
d^{\mathrm{ell}}([a_m | a_n] \otimes \eta_{12}^{\mathrm{ell}})
= -\frac{\kappa_{\mathrm{ch}}\pi^2}{3}\,E_2(\tau)
\cdot z_2^{-m-n-2}\,dz_2.
\end{equation}
\end{computation}

\begin{computation}[Degree-2 curvature]\label{comp:ell-curvature}
\index{curvature!elliptic bar complex}
The elliptic bar complex of $\cH_{\kappa_{\mathrm{ch}}}$ has curvature
element $m_0^{\mathrm{ell}} \in \B^{\mathrm{ell}}_0(\cH_{\kappa_{\mathrm{ch}}})$.
From the genus-$1$ propagator
(Vol~I, Theorem~\ref{thm:heisenberg-genus-one-complete}), the
two-point function on $E_\tau$ is $\langle a(z_1)\,a(z_2)
\rangle_{E_\tau} = \kappa_{\mathrm{ch}}\,G_\tau(z_1 - z_2)$, where
$G_\tau(z) = \zeta_\tau(z) - (\pi^2 E_2(\tau)/3)\,z$.
The curvature arises from the failure of double periodicity
(cf.\ the proof of
Vol~I, Theorem~\ref{thm:heisenberg-genus2-obstruction}):
\begin{equation}\label{eq:ell-curvature}
m_0^{\mathrm{ell}}
= \kappa_{\mathrm{ch}} \cdot \frac{\pi^2 E_2(\tau)}{3}.
\end{equation}
Expanding in the nome $q = e^{2\pi i \tau}$:
\begin{equation}\label{eq:ell-curvature-q}
m_0^{\mathrm{ell}}
= \frac{\kappa_{\mathrm{ch}}\pi^2}{3}\bigl(
1 - 24q - 72q^2 - 96q^3 - \cdots
\bigr).
\end{equation}
The sign convention: we take $m_0^{\mathrm{ell}}$ to be the generator of degree-$0$ curvature in the $\Eone$-chiral bar complex, related to the degree-$2$ bar differential~\eqref{eq:ell-bar-diff-result} by the standard curved-$A_\infty$ identity $d_{\mathrm{bar}}|_{\mathrm{deg}\,2} = -m_0 \cdot (\text{vacuum pairing})$; the sign difference between~\eqref{eq:ell-bar-diff-result} and~\eqref{eq:ell-curvature} is that identity.
At the rational limit $\tau \to i\infty$ ($q \to 0$),
$E_2(\tau) \to 1$ and $m_0^{\mathrm{ell}} \to
\kappa_{\mathrm{ch}}\pi^2/3$, recovering the rational curvature
(up to the normalization convention for the propagator
on $\bC$ versus $E_\tau$; on $\bC$ the propagator
$1/z$ has no $B$-cycle and hence $m_0^{\mathrm{rat}} = 0$).
\end{computation}

\subsection{Comparison table: elliptic versus rational}
\label{subsec:ell-vs-rat-table}

\begin{table}[ht]
\centering
\caption{Laurent coefficients controlling elliptic corrections
for $\cH_{\kappa_{\mathrm{ch}}}$}
\label{tab:ell-vs-rat-bar}
\begin{tabular}{|c|c|c|c|}
\hline
\textbf{Correction order} & \textbf{Rational coefficient}
& \textbf{Elliptic coefficient}
& \textbf{Modular type} \\
\hline
$0$ & $1/\epsilon$ & $1/\epsilon$ & polar part \\
\hline
$1$ & $0$ & $-\frac{\pi^2}{3}\,E_2(\tau)\epsilon$
& $E_2(\tau)$ (quasi-modular, wt~$2$) \\
\hline
$2$ & $0$ & $-\frac{\pi^4}{45}\,E_4(\tau)\epsilon^3$
& $E_4(\tau)$ (modular, wt~$4$) \\
\hline
$3$ & $0$ & $-\frac{2\pi^6}{945}\,E_6(\tau)\epsilon^5$
& $E_6(\tau)$ (modular, wt~$6$) \\
\hline
$4$ & $0$ & $-\frac{\pi^8}{4725}\,E_8(\tau)\epsilon^7$
& $E_8(\tau)$ (modular, wt~$8$) \\
\hline
$5$ & $0$ & $-\frac{2\pi^{10}}{93555}\,E_{10}(\tau)\epsilon^9$
& $E_{10}(\tau)$ (modular, wt~$10$) \\
\hline
$6$ & $0$ & $-\frac{1382\pi^{12}}{638512875}\,E_{12}(\tau)\epsilon^{11}$
& $E_{12}(\tau)$ (modular, wt~$12$) \\
\hline
$7$ & $0$ & $-\frac{4\pi^{14}}{18243225}\,E_{14}(\tau)\epsilon^{13}$
& $E_{14}(\tau)$ (modular, wt~$14$) \\
\hline
\textbf{General $k$}
& $0$ & $c_k(\tau)\epsilon^{2k-1}$
& weight $2k$ \\
\hline
\end{tabular}
\end{table}

\noindent
The rational propagator has only the polar coefficient $1/\epsilon$.
On $E_\tau$, the $B$-cycle monodromy of $\zeta_\tau$ introduces the
regular Eisenstein coefficients $c_k(\tau)$. For the bare Heisenberg
current OPE $a(z)a(w)\sim\kappa_{\mathrm{ch}}(z-w)^{-2}$, a single
adjacent contraction sees only $c_1(\tau)$, as in
Equation~\eqref{eq:ell-bar-diff-result}. The higher $c_k(\tau)$ enter
through descendant insertions, higher-pole fields, or iterated
correction terms in the bar filtration.

\begin{proposition}[Filtration of the elliptic bar complex]\label{prop:ell-bar-decomposition}
\index{bar complex!elliptic!decomposition}
The elliptic bar complex of $\cH_{\kappa_{\mathrm{ch}}}$ at degree $n$
carries a finite correction-order filtration
\[
0=\mathcal F^n\B^{\mathrm{ell}}_n
\subset \mathcal F^{n-1}\B^{\mathrm{ell}}_n
\subset\cdots\subset \mathcal F^1\B^{\mathrm{ell}}_n
\subset \mathcal F^0\B^{\mathrm{ell}}_n
=\B^{\mathrm{ell}}_n(\cH_{\kappa_{\mathrm{ch}}})
\]
with $\operatorname{gr}^0_{\mathcal F}\B^{\mathrm{ell}}_n
\cong \B^{\mathrm{rat}}_n(\cH_{\kappa_{\mathrm{ch}}})$. The quotient
$\operatorname{gr}^k_{\mathcal F}\B^{\mathrm{ell}}_n$ is generated by
terms with total Eisenstein correction weight~$k$, hence has
coefficients in quasi-modular forms of weight~$2k$. In particular
$\operatorname{gr}^k_{\mathcal F}\B^{\mathrm{ell}}_n=0$ for $k\geq n$.
\end{proposition}

\begin{proof}
A degree-$n$ bar element $[a_{m_1} | \cdots | a_{m_n}]$ carries
$\binom{n}{2}$ pair-Wick factors $\eta_{ij}^{\mathrm{ell}}$ indexed
by $1 \leq i < j \leq n$ (the full configuration-space propagator
pattern). The bar differential acts by adjacent contraction, so
each application of $d$ involves $n - 1$ adjacent slots
$(i, i{+}1)$, and the correction-order filtration is
driven by these adjacent slots: $k$ Eisenstein corrections require
$k$ distinct adjacent slots, bounded above by $n - 1$.

Each adjacent propagator is expanded via~\eqref{eq:zeta-laurent} as
$\zeta_\tau(\epsilon_{i,i+1}) = \epsilon_{i,i+1}^{-1} + \sum_{k \geq 1}
c_k(\tau)\,\epsilon_{i,i+1}^{2k-1}$. The bar differential extracts
the residue at the diagonal $\epsilon_{i,i+1} = 0$; since the OPE
$a(z_1)\,a(z_2) \sim \kappa_{\mathrm{ch}}/(z_1 - z_2)^2$ has a
double pole, the residue at each adjacent slot involves at most
the $\epsilon^{-1}$ coefficient of the integrand. The Laurent
term $c_j(\tau)\epsilon^{2j-1}$ carries correction weight~$j$. The
single current-current contraction contributes only the weight-$1$
coefficient $c_1(\tau)$, which gives
Equation~\eqref{eq:ell-residue-expansion}. Products of correction terms
from several adjacent slots have total weight equal to the sum of their
individual weights, and the quasi-modular weight is twice this number.
There are at most $n-1$ adjacent slots, so no term of correction order
$k\geq n$ can occur.
\end{proof}


\section{\texorpdfstring{Dynamical Yang--Baxter equation for $\mathfrak{sl}_2$}{Dynamical Yang--Baxter equation for sl2}}
\label{sec:dynamical-ybe-sl2}
\index{dynamical Yang--Baxter equation|textbf}
\index{Yang--Baxter equation!dynamical}

What relation on a meromorphic $R(u,\lambda)$ stands in for $d^2 = 0$ when the propagator lives on $E_\tau$? The ordinary YBE controls consistency of unordered double contractions in the rational bar complex; at genus~$1$ the naive Arnold relation fails (Theorem~\ref{thm:elliptic-arnold-fails}), the $R$-matrix is forced to depend on a second parameter $\lambda$, and the coherence of the three-fold contraction with respect to both parameters is the dynamical YBE. The reduction we establish below is: DYBE at $(u,v,u{+}v) \in E_\tau^3$ is the quartic Fay trisecant identity (Theorem~\ref{thm:fay}) written in $R$-matrix form.

Throughout this section we fix the base Lie algebra and representation: $\mathfrak{g} = \mathfrak{sl}_2$, with Cartan generator $h \in \mathfrak{h}$ acting on $V = \bC^2$ as $h = \mathrm{diag}(+1,-1)$ in the weight basis $\{v_+, v_-\}$ (so $V_{\pm 1}$ have weight $\pm 1$ under $h$). On $V^{\otimes 3}$ we write $h^{(i)}$ for the endomorphism acting as $h$ on the $i$-th tensor factor and as identity elsewhere; on each basis vector $v_{\epsilon_1}\otimes v_{\epsilon_2}\otimes v_{\epsilon_3}$ the operator $h^{(i)}$ acts by the scalar $\epsilon_i \in \{+1,-1\}$. The dynamical parameter $\lambda$ lives in $\mathfrak{h}^* \cong \bC$, dual to $h$.

\subsection{The dynamical Yang--Baxter equation}
\label{subsec:dybe-statement}

\begin{definition}[Dynamical Yang--Baxter equation]
\label{def:dybe}
\index{dynamical Yang--Baxter equation!definition}
With $\mathfrak{g} = \mathfrak{sl}_2$, $V = \bC^2$, $h \in \mathfrak{h}$ and $h^{(i)}$ as above, the \emph{dynamical Yang--Baxter equation} (DYBE) for an
$R$-matrix $R(u, \lambda) \in \mathrm{End}(V \otimes V)$
depending on spectral parameter $u$ and dynamical parameter
$\lambda \in \mathfrak{h}^*$ is:
\begin{equation}\label{eq:dybe}
R_{12}(u, \lambda - \gamma h^{(3)})\,
R_{13}(u+v, \lambda)\,
R_{23}(v, \lambda - \gamma h^{(1)})
\;=\;
R_{23}(v, \lambda)\,
R_{13}(u+v, \lambda - \gamma h^{(2)})\,
R_{12}(u, \lambda)
\end{equation}
where $h^{(i)}$ denotes the action of $h \in \mathfrak{h}$ on
the $i$-th tensor factor of $V^{\otimes 3}$, and
$\lambda - \gamma h^{(i)}$ means that the dynamical parameter
$\lambda$ is shifted by $\gamma$ times the weight of the state in
factor~$i$ ($\gamma$ the Felder quantization parameter; conventions
as in Felder--Varchenko~\cite{Felder94,FelderVarchenko96}).
\end{definition}

\subsection{\texorpdfstring{Explicit verification for $\mathfrak{sl}_2$}{Explicit verification for sl2}}
\label{subsec:dybe-sl2-verification}

For $\mathfrak{sl}_2$ with $V = \bC^2$, the weight decomposition
is $V = V_+ \oplus V_-$ with $h \cdot v_\pm = \pm v_\pm$.
We use the Felder elliptic $R$-matrix of Definition~\ref{def:elliptic-r}:
\begin{equation}\label{eq:sl2-r-matrix-entries}
R(u, \lambda) =
\begin{pmatrix}
1 & 0 & 0 & 0 \\
0 & \alpha(u, \lambda) & \beta(u, \lambda) & 0 \\
0 & \beta(u, -\lambda) & \alpha(u, -\lambda) & 0 \\
0 & 0 & 0 & 1
\end{pmatrix}
\end{equation}
with $\alpha(u,\lambda) = \theta(u)\,\theta(\lambda+\gamma)/\allowbreak
(\theta(u-\gamma)\,\theta(\lambda))$ and
$\beta(u,\lambda) = -\theta(u+\lambda)\,\theta(\gamma)/\allowbreak
(\theta(u-\gamma)\,\theta(\lambda))$
(the entries from Definition~\ref{def:elliptic-r},
abbreviated here as $\theta = \theta_1(\cdot\,|\,\tau)$; $\gamma$ is the Felder quantization parameter per Remark~\ref{rem:eta-to-gamma}).

\begin{computation}[Matrix entries of the DYBE]
\label{comp:dybe-matrix-entries}
We examine the DYBE~\eqref{eq:dybe} on the weight spaces of
$V^{\otimes 3}$. On the space $V^{\otimes 3}$
with basis $\{v_{\epsilon_1} \otimes v_{\epsilon_2}
\otimes v_{\epsilon_3} : \epsilon_i \in \{+,-\}\}$
(eight basis vectors), the dynamical shifts act as
$h^{(i)} \cdot v_{\epsilon_1\epsilon_2\epsilon_3}
= \epsilon_i\,v_{\epsilon_1\epsilon_2\epsilon_3}$.
The DYBE is a system of $64$ scalar equations (the entries of
an $8 \times 8$ matrix equation).

\emph{Extremal entries.}
On the extremal weight vectors $v_{+++}$ and $v_{---}$ every factor
acts by its unit corner entry, so both sides of~\eqref{eq:dybe} equal
$1$: the equation is trivially satisfied.

\emph{Mixed weight spaces.}
The nontrivial content lies on the $3$-dimensional subspaces
spanned by $\{v_{++-}, v_{+-+}, v_{-++}\}$ (weight~$1$)
and $\{v_{--+}, v_{-+-}, v_{+--}\}$ (weight~$-1$).
By the symmetry $R(u,\lambda) = R(u,-\lambda)|_{+\leftrightarrow -}$,
it suffices to verify the weight-$1$ block, a $3 \times 3$ matrix
equation.

\emph{Diagonal entries of the weight-$1$ block.}
On $v_{++-}$, the coefficient of $v_{++-}$ on the left
of~\eqref{eq:dybe} is
$\alpha(v,\lambda{-}\gamma)\,\alpha(u{+}v,\lambda)$ and on the right
$\alpha(u{+}v,\lambda{-}\gamma)\,\alpha(v,\lambda)$; the two agree
because $\alpha(z,\mu) =
\bigl[\theta(z)/\theta(z{-}\gamma)\bigr]\cdot
\bigl[\theta(\mu{+}\gamma)/\theta(\mu)\bigr]$
factorises the $z$- and $\mu$-dependence.

\emph{Off-diagonal entries.}
These are genuine theta identities. The coefficient of $v_{+-+}$
in~\eqref{eq:dybe} applied to $v_{++-}$ reads
\begin{equation}\label{eq:dybe-offdiag-entry}
\alpha(v,\lambda{-}\gamma)\,\beta(u{+}v,-\lambda)\,
\beta(u,\lambda{-}\gamma)
+ \beta(v,\gamma{-}\lambda)\,\alpha(u,\lambda{-}\gamma)
\;=\;
\alpha(u{+}v,\lambda{-}\gamma)\,\beta(v,-\lambda).
\end{equation}
Substituting the entries, multiplying through by
$\theta(v{-}\gamma)\,\theta(\lambda{-}\gamma)^2\,\theta(u{-}\gamma)\,
\theta(u{+}v{-}\gamma)/\theta(\gamma)$ and writing
$[x] := \theta(x|\tau)$, equation~\eqref{eq:dybe-offdiag-entry}
becomes the quartic three-term identity
\begin{equation}\label{eq:dybe-offdiag-quartic}
[\lambda]\,[u]\,[v{-}\lambda{+}\gamma]\,[u{+}v{-}\gamma]
\;-\; [\gamma]\,[v]\,[u{+}\lambda{-}\gamma]\,[u{+}v{-}\lambda]
\;=\; [\lambda{-}\gamma]\,[u{-}\gamma]\,[u{+}v]\,[v{-}\lambda],
\end{equation}
the cleared kernel form~\eqref{eq:fay-quartic-cleared} of the Fay
identity at $(z_1, z_2, \lambda_1, \lambda_2) =
(u{+}v,\, v,\, -\gamma,\, \gamma{-}\lambda)$.
\end{computation}

\begin{theorem}[Felder; DYBE reduces to Fay]
\label{prop:dybe-reduces-to-fay}
\index{Fay trisecant identity!and dynamical YBE}
The $\mathfrak{sl}_2$ elliptic $R$-matrix of
Definition~\ref{def:elliptic-r} satisfies the dynamical Yang--Baxter
equation~\eqref{eq:dybe}. Each scalar entry of~\eqref{eq:dybe}
reduces, after clearing denominators, either to the factorisation
identity of the diagonal entries or to an instance of the quartic Fay
trisecant identity in the cleared kernel
form~\eqref{eq:fay-quartic-cleared}
\textup{(}Theorem~\textup{\ref{thm:fay})}, together with the
quasi-periodicity $\theta(z+1|\tau) = -\theta(z|\tau)$,
$\theta(z+\tau|\tau) = -e^{-\pi i\tau - 2\pi i z}\theta(z|\tau)$.
\end{theorem}

\begin{proof}
The extremal entries and the diagonal entries of the weight-$1$ block
are settled in Computation~\ref{comp:dybe-matrix-entries}. The
displayed off-diagonal entry~\eqref{eq:dybe-offdiag-entry} reduces to
the Fay identity~\eqref{eq:fay-quartic-cleared} at
$(z_1, z_2, \lambda_1, \lambda_2) = (u{+}v, v, -\gamma,
\gamma{-}\lambda)$ (Computation~\ref{comp:dybe-matrix-entries}). The
remaining off-diagonal entries of the weight-$1$ block reduce
to~\eqref{eq:fay-quartic-cleared} at permuted specializations by the
same clearing procedure, and the weight-$(-1)$ block follows from the
weight-$1$ block by the symmetry
$R(u,\lambda) = R(u,-\lambda)|_{+\leftrightarrow -}$; the systematic
entry-by-entry verification is
Felder--Varchenko~\cite{Felder94,FelderVarchenko96}. Quasi-periodicity
enters when matching entries whose $\lambda$-arguments differ by
lattice translates.
\end{proof}

\begin{conjecture}[Dynamical elliptic bar nilpotency]
\label{conj:dynamical-bar-nilpotency}
\ClaimStatusConjectured
\index{bar complex!elliptic!nilpotency via DYBE}
There is a $\lambda$-extended bar complex
$\B^{\mathrm{ell}}_\lambda(E_{q,p}(\mathfrak{sl}_2))$ of the elliptic
quantum group
\textup{(}Definition~\textup{\ref{def:elliptic-quantum})} --- bar
words tensored with functions of the dynamical parameter, the
differential contracting with the Felder $R$-matrix and shifting
$\lambda \mapsto \lambda - \gamma h^{(i)}$ --- such that
$d_\lambda^2 = 0$ on $\overline{C}_3(E_\tau)$, with the dynamical
Yang--Baxter equation~\eqref{eq:dybe} as the precise coherence law:
$d_\lambda^2 = 0$ if and only if $R(u,\lambda)$ satisfies the DYBE.
\end{conjecture}

\begin{remark}[Status and mechanism]
\label{rem:dynamical-bar-status}
Three parts of the intended equivalence are theorems; the remaining
obligation is the construction. \emph{Proved:} the Felder $R$-matrix
satisfies the DYBE (Theorem~\ref{prop:dybe-reduces-to-fay}); the
naive elliptic bar differential, without the dynamical extension,
does \emph{not} square to zero --- the nine double-contraction terms
on $\overline{C}_3(E_\tau)$ (three from each of the ordered
contractions $(12)(23)$, $(23)(31)$, $(31)(12)$) no longer cancel,
their obstruction being the quasi-modular defect
$S_\tau\,da\wedge db$ of Theorem~\ref{thm:elliptic-arnold-fails};
and the $B$-cycle monodromy
$\zeta_\tau(z+\tau) = \zeta_\tau(z) - 2\pi i$ generating that defect
is the monodromy tracked by the dynamical shift
$\lambda \mapsto \lambda - \gamma h^{(i)}$
(Corollary~\ref{cor:elliptic-bar-dynamical};
cf.\ Remark~\ref{rem:elliptic-qg-bar}(ii)). \emph{Open:} the
construction of the $\lambda$-extended complex itself --- the correct
completion of bar words against functions of $\lambda$, the
identification of the shifted differential with the KZB connection,
and the verification that the grouped nine-term obstruction is
exactly the DYBE deviation. The kernel-form Fay
identity~\eqref{eq:fay-kernel}, whose second argument is the
dynamical variable, is the identity from which the cancellation must
come.
\end{remark}


\section{Shuffle algebra presentation}
\label{sec:shuffle-algebra}
\index{shuffle algebra|textbf}
\index{toroidal algebra!shuffle presentation}

The Drinfeld presentation writes $U_{q,t}(\hat{\hat{\fg}})$ by generators and
relations, with the relations encoded in the elliptic $R$-matrix. Computing
inside this presentation means reducing words modulo Serre and RLL, which
is intractable beyond the lowest modes. The shuffle algebra linearises the
problem: it replaces the generators-and-relations presentation by a
function-theoretic one on the symmetric algebra
$\bigoplus_n \bC(z_1,\ldots,z_n)^{S_n}$, with the elliptic structure
function $\xi(u) = \theta(u+\gamma)/\theta(u)$ absorbed into the product
rule (the symbol $\xi$ disambiguates the shuffle structure function from the log-derivative $\zeta_\tau$ and the Weierstrass $\zeta_W$ used in \S\ref{sec:elliptic-bar-heisenberg}; $\gamma$ is the Felder quantization parameter of Remark~\ref{rem:eta-to-gamma}). Under residue duality the shuffle product on $\cS_n$ pairs with the deshuffle coproduct on $\B_n(\cA)$ of Vol~I, so shuffle computations are bar computations transported to the function-theoretic side.

\subsection{Definition of the shuffle algebra}
\label{subsec:shuffle-def}

\begin{definition}[Shuffle algebra]\label{def:shuffle-algebra}
\index{shuffle algebra!definition}
Let $\xi(u) = \theta(u+\gamma)/\theta(u)$ be the structure function
determined by the elliptic $R$-matrix (Definition~\ref{def:elliptic-r}): in kernel form
$\xi(u) = \bigl[\theta(\gamma)/\theta'(0)\bigr] F_\tau(u,\gamma)$,
the Kronecker--Eisenstein kernel of Theorem~\ref{thm:fay} evaluated
at dynamical argument $\gamma$.
The \emph{shuffle algebra} $\mathcal{S}_\xi$ is the $\bZ_{\geq 0}$-graded
vector space
\[
\mathcal{S}_\xi = \bigoplus_{n \geq 0} \mathcal{S}_n, \quad
\mathcal{S}_n = \bC(z_1, \ldots, z_n)^{S_n}_{\mathrm{sym}}
\]
(symmetric rational functions in $n$ variables), equipped with
the \emph{shuffle product}:
for $f \in \mathcal{S}_m$ and $g \in \mathcal{S}_n$, define
$f \star g \in \mathcal{S}_{m+n}$ by
\begin{multline}\label{eq:shuffle-product}
(f \star g)(z_1, \ldots, z_{m+n})
= \frac{1}{m!\,n!} \sum_{\sigma \in S_{m+n}}
\prod_{\substack{i \leq m \\ j > m}}
\xi(z_{\sigma(i)} - z_{\sigma(j)}) \\
\cdot\, f(z_{\sigma(1)}, \ldots, z_{\sigma(m)})
\cdot g(z_{\sigma(m+1)}, \ldots, z_{\sigma(m+n)}).
\end{multline}
Equivalently, the sum runs over $(m,n)$-shuffles
$\mathrm{Sh}(m,n) \subset S_{m+n}$, with
the $1/(m!\,n!)$ factor replaced by $1$ and the
sum restricted to shuffles.
\end{definition}

\begin{remark}[Origin]\label{rem:shuffle-origin}
The shuffle algebra was introduced by
Feigin--Odesskii in the rational and
trigonometric cases and extended to the elliptic case by
Schiffmann--Vasserot~\cite{SV13} and by Negu\c{t}~\cite{Neg14}.
The rational specialization $\xi(u) = (u+\gamma)/u$
recovers the Feigin--Odesskii shuffle realization of the
Yangian; the trigonometric specialization
$\xi(u) = (1 - q e^{2\pi i u})/(1 - e^{2\pi i u})$
recovers the quantum affine algebra. We distinguish this shuffle-algebra appearance of Schiffmann--Vasserot from the cohomological Hall algebra SV theorem $H^{\mathrm{CoHA}}(\bC^3) \cong Y^+(\widehat{\mathfrak{gl}}_1)$ (positive half only; the full affine Yangian is the Drinfeld double): the shuffle algebra linearises the generators-and-relations side, while CoHA linearises the moduli-of-sheaves side.
\end{remark}

\subsection{Explicit shuffle products}
\label{subsec:shuffle-explicit}

\begin{computation}[First generators]\label{comp:shuffle-generators}
\index{shuffle algebra!explicit products}
Let $e_1 = 1 \in \mathcal{S}_1 = \bC(z)$. Writing $\xi(u) := \theta(u+\gamma)/\theta(u)$ for the elliptic structure function and applying oddness of $\theta$ to $\xi(-z_{12}) = \theta(\gamma - z_{12})/\theta(-z_{12}) = -\theta(\gamma - z_{12})/\theta(z_{12})$:
\begin{equation}\label{eq:e1e1-result}
(e_1 \star e_1)(z_1, z_2)
= \xi(z_{12}) + \xi(-z_{12})
= \frac{\theta(z_{12}+\gamma) - \theta(\gamma - z_{12})}
{\theta(z_{12})},
\end{equation}
where $z_{12} = z_1 - z_2$ and $\gamma$ is the quantization parameter (see Remark~\ref{rem:eta-to-gamma}). The triple product is
\begin{equation}\label{eq:shuffle-e1e1e1}
(e_1 \star e_1 \star e_1)(z_1, z_2, z_3)
= \sum_{\sigma \in S_3}
\xi(z_{\sigma(1)} - z_{\sigma(2)})\,
\xi(z_{\sigma(1)} - z_{\sigma(3)})\,
\xi(z_{\sigma(2)} - z_{\sigma(3)}).
\end{equation}
Associativity of the shuffle product holds for any structure function $\xi$: both bracketings of a triple product equal the sum over $(m,n,p)$-shuffles with the kernel $\prod \xi(z_i - z_j)$ over all cross pairs. The elliptic input enters not through associativity but through the wheel conditions cutting out the sub-shuffle algebra and through the DYBE coherence of the braiding (Theorem~\ref{prop:dybe-reduces-to-fay}).
\end{computation}

\begin{remark}[Shuffle--bar isomorphism]\label{rem:shuffle-bar}
\index{shuffle algebra!bar complex connection}
For a shuffle-realised positive half with the residue pairing fixed,
$\mathcal{S}_n$ identifies with the restricted dual of the corresponding
bar degree. Under this identification the shuffle product is dual to the
deshuffle coproduct on $\B(\cA)$ (Vol~I, Corollary~\ref{cor:bar-is-dgcoalg}).
\end{remark}


\section{Comparison: rational, trigonometric, elliptic}
\label{sec:rat-trig-ell-comparison}
\index{comparison!rational vs trigonometric vs elliptic|textbf}

\begin{table}[ht]
\centering
\caption{Comparison across the three regimes}
\label{tab:rat-trig-ell}
\renewcommand{\arraystretch}{1.35}
{\footnotesize
\begin{tabular}{|l|c|c|c|}
\hline
& \textbf{Rational} & \textbf{Trigonometric} & \textbf{Elliptic} \\
\hline
\textbf{Algebra}
& Yangian $Y(\fg)$
& Qtm.\ affine $U_q(\hat{\fg})$
& Toroidal $U_{q,t}(\hat{\hat{\fg}})$ \\
\hline
\textbf{Curve}
& $\bC$ (or $\bC\bP^1$)
& $\bC^*$ (or $\bC\bP^1 \!\setminus\! \{0,\infty\}$)
& $E_\tau$ \\
\hline
\textbf{Propagator}
& $\frac{dz}{z}$
& $\frac{dz}{1 - z}$
& $d\log\theta_1(z|\tau)$ \\
\hline
\textbf{$R$-matrix}
& $1 - \hbar P/u$
& $(q{-}q^{-1})/(u{-}u^{-1})$
& Felder $R(u,\lambda)$ (Def.~\ref{def:elliptic-r}) \\
\hline
\textbf{Arnold id.}
& $\eta_{12}\!\wedge\!\eta_{23} + \text{cyc} = 0$
& holds in $w = e^{2\pi i z}$; $\pi\cot$ defect $\pi^2$
& fails: defect $S_\tau$ (Thm~\ref{thm:elliptic-arnold-fails}) \\
\hline
\textbf{$d^2\!=\!0$}
& Partial fractions
& $q$-partial fractions
& DYBE + dyn.\ shift (Conj.~\ref{conj:dynamical-bar-nilpotency}) \\
\hline
\textbf{Koszul dual}
& $Y_{-\hbar}(\fg)$
& $U_{q^{-1}}(\hat{\fg})$
& $U_{q^{-1},t^{-1}}(\hat{\hat{\fg}})$ (conj.) \\
\hline
\textbf{Struct.\ fn.}
& $(u{+}\gamma)/u$
& $(1{-}qe^u)/(1{-}e^u)$
& $\theta(u{+}\gamma)/\theta(u)$ \\
\hline
\textbf{Shuffle alg.}
& Feigin--Odesskii
& F--O ($q$-def.)
& Schiffmann--Vasserot \\
\hline
\textbf{Dyn.\ par.}
& absent
& absent
& $\lambda \in \mathfrak{h}^*$ \\
\hline
\textbf{Bar curv.}
& $m_0 = 0$
& $m_0 \propto \log q$
& $m_0 \propto E_2(\tau)$ \\
\hline
\end{tabular}
}
\end{table}

\begin{remark}[Degeneration limits]\label{rem:degeneration-limits}
\index{degeneration!elliptic to trigonometric to rational}
The three regimes degenerate via
\begin{equation}\label{eq:degeneration-chain}
\text{Elliptic} \;\xrightarrow{\;\tau \to i\infty\;}
\text{Trigonometric} \;\xrightarrow{\;q \to 1\;}
\text{Rational}.
\end{equation}
As $\tau \to i\infty$, $\theta_1(z|\tau) \to 2\sin(\pi z)$ and the Eisenstein corrections collapse to constants; as $q \to 1$, the trigonometric propagator becomes $dz/z$ and $R_q(u) \to 1 - \hbar P/u$. The bar complex degenerates accordingly: $\B^{\mathrm{ell}}(\cA) \to \B^{\mathrm{trig}}(\cA) \to \B^{\mathrm{rat}}(\cA)$, consistent with the spectral sequence of Theorem~\ref{thm:elliptic-vs-rational}.
\end{remark}


%% BRIDGE 1: E as shared geometry

\subsection{The elliptic curve as shared geometry}
\label{subsec:bridge-shared-E}
\index{elliptic curve!shared geometry bridge}

\begin{remark}[Bridge: algebras \emph{on} \texorpdfstring{$E_\tau$}{E-tau} and algebras \emph{from} \texorpdfstring{$K3 \times E$}{K3 x E}, heuristic]
\label{rem:bridge-shared-E}
The elliptic curve $E_\tau$ enters this chapter in two capacities. First,
it is the curve on which the propagator $d\log\theta_1(z|\tau)$, the
Eisenstein corrections to the bar differential
(Theorem~\ref{thm:elliptic-vs-rational}), the toroidal algebra
$U_{q,t}(\hat{\hat{\mathfrak{g}}})$ with
$\tau = \log t / \log q$ (Conjecture~\ref{conj:toroidal-e1}), and
Felder's elliptic $R$-matrix $R(u,\lambda)$
(Definition~\ref{def:elliptic-r}) live.

Second, $E_\tau$ is the elliptic factor in the Calabi--Yau threefold
$K3 \times E$. The canonical stage-$1$
$\PhiFA_3(D^b(\Coh(K3 \times E))) \in \EdHolFA(K3 \times E)$ of
Theorem~\ref{thm:phi-two-stage-factorisation} lives on the threefold
and factorises over $3$-dimensional opens. If the Costello--Li
holomorphic twist reduction along $K3$ realises the specialisation
$\SpCh_{\Sigma_2^{K3}, E_\tau}$ with
$\Sigma_2^{K3} = K3 \times \{\mathrm{pt}\}$ and reference curve
$C = E_\tau$, it descends this stage-$1$ factorisation algebra to a
boundary $\Eone$-chiral algebra $A_E$ on $E_\tau$ (Construction~\ref{constr:costello-li-k3xe}).
On the rank-$24$ Heisenberg hypothesis, $c(A_E)=24$ is the Mukai-lattice
rank $\kappa_{\mathrm{fiber}}=24$, not the compact total-space
$\kappa_{\mathrm{cat}}(K3\times E)=0$ and not the default
$K3\times E$ chiral value $\kappa_{\mathrm{ch}}^{\mathrm{Heis}}=3$.
The BKM-side value is separate:
$\kappa_{\mathrm{BKM}}(\Phi_N)=c_N(0)/2$; in the default $N=1$ denominator
normalisation for $\Delta_5$, $c_1(0)=10$ and
$\kappa_{\mathrm{BKM}}(\Delta_5)=5$. Thus the
$K3\times E$ spectrum is $\{0, 0, 3, 5, 24\}$ from five distinct
constructions: $\kappa_{\mathrm{cat}} = 0$ (Künneth multiplicative
on the total space), $\kappa_{\mathrm{ch}}^{\mathrm{Hodge}} = 0$
(compact total-space Hodge supertrace, Serre at odd $d = 3$),
$\kappa_{\mathrm{ch}}^{\mathrm{Heis}} = 3$ (Stage-$2$ Heisenberg
rank), $\kappa_{\mathrm{BKM}}(\Delta_5) = 5$ (Borcherds weight at
$N = 1$), $\kappa_{\mathrm{fiber}} = 24$ (K3 Mukai-lattice rank). The existence of $A_E$ and the Heisenberg reduction are
conjectural at $d=3$; the identification is a heuristic bridge.

If the bridge holds, the toroidal algebras
$U_{q,t}(\hat{\hat{\mathfrak{g}}})$, Felder's elliptic
$E_{q,p}(\mathfrak{g})$, and the conjectural Costello--Li boundary
$A_E$ are chiral algebras on the same curve $E_\tau$, sharing the
propagator $d\log\theta_1$, the $E_2$-curved differential with its
dynamical correction (Theorem~\ref{thm:elliptic-arnold-fails},
Corollary~\ref{cor:elliptic-bar-dynamical}), and the Eisenstein
filtration of the bar complex. They differ in the target algebra:
$A_E$ would have $24$ free-boson generators and shadow class~G
($r_{\max}=2$); $E_{q,p}(\mathfrak{sl}_N)$ has $N^2-1$ generators and
a dynamical parameter $\lambda \in \mathfrak{h}^*$.
\end{remark}

\subsection{Summary}
\label{subsec:toroidal-elliptic-summary}

Three curve-based structures on $E_\tau$ share one differential. The
propagator $d\log\theta_1(z|\tau)$, the quartic Fay trisecant
identity with the quasi-modular defect it governs
(Theorems~\ref{thm:fay} and~\ref{thm:elliptic-arnold-fails}), and the
Eisenstein filtration of the bar complex assemble into a single
$\Eone$-chiral package. The quantum toroidal
$U_{q,t}(\hat{\hat{\mathfrak{g}}})$ of
Conjecture~\ref{conj:toroidal-e1}, the Felder dynamical
$E_{q,p}(\mathfrak{g})$ of Theorem~\ref{prop:dybe-reduces-to-fay}
and Conjecture~\ref{conj:dynamical-bar-nilpotency}, and the elliptic
shuffle algebra of Definition~\ref{def:shuffle-algebra} are three
presentations of this curve-level package at their stated scopes. The Koszul-dual pairing
$(q,t) \leftrightarrow (q^{-1},t^{-1})$ inverts orientation on
$E_\tau$ via $\tau \mapsto -\tau$, the elliptic shadow of the Verdier
chiral duality $\mathbb{D}_{\mathrm{Ran}}\bar{B}(A) \simeq \bar{B}(A^!)$.

If the specialisation $\SpCh_{\Sigma_2,E_\tau}$ of the stage-$1$
$\PhiFA_3$ on $K3 \times E$ lands in this $\Eone$-chiral package with a
rank-$24$ Mukai Heisenberg boundary target, then the boundary algebra
$A_E$ inherits an $E_2$-curved bar differential
(Corollary~\ref{cor:elliptic-bar-dynamical}(i)) with central charge
$24$. That scalar is the fibre/Mukai datum $\kappa_{\mathrm{fiber}}$;
the compact $K3\times E$ values remain
$\kappa_{\mathrm{cat}} = 0$, $\kappa_{\mathrm{ch}}^{\mathrm{Hodge}} = 0$,
$\kappa_{\mathrm{ch}}^{\mathrm{Heis}} = 3$,
$\kappa_{\mathrm{BKM}}(\Delta_5) = 5$, and $\kappa_{\mathrm{fiber}} = 24$.



\section{Hall positive half on \texorpdfstring{$K3 \times E$}{K3 x E} and the elliptic chart}
\label{sec:tor-ell-k3-lift}

\begin{remark}[Hall positive half and centre-side braiding over \texorpdfstring{$K3 \times E$}{K3 x E}]
\label{rem:tor-ell-k3-lift-opener}
The comparison with $K3\times E$ has three layers. First, the framed
CY$_3$ specialisation
\[
A_{\Sigma_2,C}
=\SpCh_{\Sigma_2,C}\bigl(\PhiFA_3(D^b(\Coh(K3\times E)))\bigr)
\]
is an $\Eone$-chiral algebra on the reference curve $C$ whenever the
hypotheses of Theorem~\ref{thm:cy-to-chiral-d3} are met. The native
stage-$1$ object is the holomorphic factorisation algebra on the
threefold; the $\Etwo$ structure relevant to the $R$-matrix is recovered
on the centre side $\cZ(\Rep^{\Eone}(A_{\Sigma_2,C}))$.

Second, the Hall comparison is a positive-half statement. The available
target is a conjectural oriented critical CoHA positive half
$Y^+_{\Hall}(K3\times E)$ whose graded character is constrained by the
reduced DT/Igusa denominator. Schiffmann--Vasserot shuffle methods and
Davison's CY$_3$ CoHA technology provide the local models for this
positive half; they do not by themselves construct the compact
$K3\times E$ double.

Third, the algebra-level endpoint
\[
G(K3\times E)=D\bigl(Y^+_{\Hall}(K3\times E)\bigr)
\]
requires a nondegenerate Hopf pairing, negative and Cartan halves,
topological completions, bracket compatibility with the
Borcherds--Kac--Moody denominator algebra, and compatibility of the
resulting braiding with $\cZ(\Rep^{\Eone})$. The K3 Mukai
self-mirror Yangian and quantum-toroidal branch supply stable-envelope
evidence for the centre-side braiding, but they are not the compact
$K3\times E$ Hall--Borcherds double. The Miki generators of the
toroidal branch therefore remain chart-level evidence until these data
are constructed.
\end{remark}
